% Landau's ceiling step on windows, with explicit constants.
% Phase 1 deliverable of the landau-windows project (see ../README.md, ../docs/).
% Compiles with plain LaTeX + AMS classes; no exotic packages.
%
% STATUS: working single-zero certificate draft; the full arithmetic
% Landau reduction is not yet proved. Absolute constants are listed in
% the Constant Registry (Section 9); derived real quantities need certified
% rational enclosures. The quoted standard estimates (Inputs 5.2--5.4:
% low zero data; local zero count; the zeta'/zeta partial-fraction bound)
% are cited with deliberately conservative constants and are flagged for
% tightening / self-containment during the effective-shadow work.

\documentclass[11pt]{amsart}

\usepackage{amsmath,amssymb,amsthm}
\usepackage[T1]{fontenc}
\usepackage{lmodern}
\usepackage{longtable}
\usepackage[colorlinks=true,linkcolor=blue]{hyperref}

\newtheorem{theorem}{Theorem}[section]
\newtheorem{lemma}[theorem]{Lemma}
\newtheorem{corollary}[theorem]{Corollary}
\newtheorem{proposition}[theorem]{Proposition}
\newtheorem{definition}[theorem]{Definition}
\newtheorem{remark}[theorem]{Remark}
\newtheorem{assumption}[theorem]{Input}

\numberwithin{equation}{section}

% Notation macros
\newcommand{\R}{\mathbb{R}}
\newcommand{\C}{\mathbb{C}}
\newcommand{\Q}{\mathbb{Q}}
\newcommand{\Z}{\mathbb{Z}}
\newcommand{\ZZ}{\zeta}
\newcommand{\res}{\operatorname{Res}}
\newcommand{\re}{\operatorname{Re}}
\newcommand{\im}{\operatorname{Im}}
\newcommand{\Sig}{\sigma}
\newcommand{\Lam}{\Lambda}
\newcommand{\Euler}{\gamma}
\newcommand{\dbar}{\bar{d}}

\begin{document}

\title{Landau's ceiling step on windows, with explicit constants\\
\large I. The effective windowed Landau theorem for the Riemann zeta function}

\author{the \texttt{landau-windows} project}

\date{November 21, 2026. Working draft (Phase 1).}

\begin{abstract}
This draft develops an explicit single-zero certificate version of the
``windowed Landau theorem'': if the smoothed prime-power counting error $G(x)$ satisfies
the one-sided bound $G(x) \le C x^{\theta}$ on a single long window
$[Y, Y^{3/2}]$, then a certified zero of the Riemann zeta function with
real part $>\theta$ satisfying the fixed-slab isolation conditions cannot
exist. Those conditions restrict clusters and multiplicities; removing
that restriction is required for the full arithmetic reduction. In the
contrapositive direction, from a zero-isolation
certificate and the ceiling data we compute an explicit \emph{address}
$x \in [Y, Y^{3/2}]$ at which the ceiling is violated, together with an
explicit margin; the verification at $x$ is a finite sum over prime powers
checked by rational interval arithmetic. The displayed absolute constants
are rational; derived real quantities require certified enclosures.
The non-effective statement we
effectivize is due to Rubinstein (2026), building on Landau (1905) and the
P\'olya--Ingham high-derivative method; our contribution is the
explicitness, the certificate formulation, and the computable address. This
is the classical layer (target T1) of a program whose ultimate aim is an
ACL2 proof of this step with a documented interpretation in weak arithmetic,
targeting approximately Peano-arithmetic strength. The finite evaluators,
analytic estimates and final composition are still being verified.
\end{abstract}

\maketitle

% Short running heads (the full title is too long for the header).
\pagestyle{myheadings}
\markboth{Landau's ceiling step on windows, with explicit constants}%
        {Landau's ceiling step on windows, with explicit constants}

\tableofcontents

\section{Introduction}\label{sec:intro}

\subsection{The step this paper effectivizes}

The Riemann hypothesis has a three-step architecture (we follow the
explanatory video guide of \cite{RieVideo} and the explanatory draft of
Rubinstein \cite{Rub26}, which in turn explain the announced Lean
development of Alp\"oge \cite{Alp26}):

\begin{enumerate}
\item zeros of the zeta function mirror across the critical line, so it
suffices to exclude zeros \emph{on the right} of the line;
\item \textbf{Landau's ceiling step}: a \emph{one-sided} bound on the
prime-counting error excludes zeros on the right --- even when the bound is
available only on recurring windows $[Y, Y^{3/2}]$;
\item building the ceiling (the arithmetic core).
\end{enumerate}

This repo is occupied with step 2, and specifically with the fact that
the windowed form of step 2 has a one-page non-effective argument
\cite{Rub26} (Part~II, \S1). The present draft studies explicit certificate
data and constants for a restricted single-zero version. The repository
separately proves the classical windowed statement in Lean and is developing
the full arithmetic reduction in ACL2.

To state what ``the ceiling'' is, let $\Lam$ denote the von Mangoldt
function and define, following the Alp\"oge development's conventions
(\texttt{Sz}, \texttt{Gz} in \texttt{StatementOC.lean}; see
\texttt{../docs/statement.md}),
\begin{equation}\label{eq:SG}
S(x) \;=\; \sum_{n \le x} \Lam(n)\Bigl(\frac{x}{n} - 1\Bigr),
\qquad
G(x) \;=\; S(x) - x\log x + (1+\Euler)\,x,
\end{equation}
where $\Euler$ is the Euler--Mascheroni constant. The statistic $S$ is the
Chebyshev-style sum weighted by the ramp $(x/n-1)^+$, and $G$ is its
deviation from the main term; $G$ oscillates around $\log 2\pi$ with
amplitude governed by the zeros of $\zeta$ (a zero with real part $\beta$
forces fluctuations of size $x^{\beta}$).

The non-effective statement we take as our starting point is the
following, stated and proved in one page in \cite{Rub26} (Part~I, \S1.1, and
Part~II, \S1).

\begin{lemma}[Rubinstein, ``the same, on intervals'']\label{lem:rubinstein}
Let $0 < \theta < 1$. If there are arbitrarily large $Y$ such that
$G(x) \le x^{\theta}$ for all $x \in [Y, Y^{3/2}]$, then
$\zeta(s) \ne 0$ for $\re s > \theta$.
\end{lemma}

Its one-page proof replaces Landau's theorem on Mellin transforms of
one-signed functions \cite{Lan05, MV07} by an estimate of high derivatives
of the Mellin transform $F$ of $G$: a high derivative ``sees only one
interval at a time'', so a ceiling that holds on a single window $[Y,
Y^{3/2}]$ already forces the high derivatives at a suitable scale to be
small, while a nearby pole of $F$ (i.e.\ a nearby zero of $\zeta$) forces
them to be large. The kernel of the argument is the P\'olya--Ingham
high-derivative method \cite{Ing32}.

\subsection{What is proved here}

Lemma~\ref{lem:rubinstein} is an existence proof by contradiction: it
concludes that a ceiling and a rogue zero cannot both stand, but it
neither says \emph{where} the ceiling breaks, nor \emph{by how much}, nor
with \emph{which} constants. The theorem of this paper (Theorem~\ref{thm:main})
is the effective refinement. Informally:

\begin{quote}
\textbf{Theorem (informal).} Fix $\tfrac12 \le \theta < 1$ and ceiling
constants $C \ge 1$. Given (i) a zero-isolation certificate: a rational
rectangle $R$ in $\re s > \theta$ certified to isolate exactly one simple
zero $\rho$ of $\zeta$, inside a certified zero-free slab to the right of
the critical line; and (ii) a window $[Y, Y^{3/2}]$ with $Y$ at least an
explicitly computable threshold $Y_0$ --- one can \emph{compute} an address
$x \in [Y, Y^{3/2}]$ and a margin $m > 0$ such that
\[
G(x) \;\ge\; (C + m)\, x^{\theta},
\]
and the verification of this inequality at $x$ is a finite sum over prime
powers compared with rational interval arithmetic.
\end{quote}

Contrapositive reading: \emph{if the ceiling claim is verified to hold
throughout a large window, then no certificate-bearing zero can exist to
the right of $\theta$}. Certificate reading: \emph{a rogue zero and a
ceiling cannot both stand; we name where the ceiling breaks, by how much,
and how to check it finitely}. This is exactly the contract required by the
surrounding architecture: the arithmetic side (step 3, not our concern)
produces ceilings on windows of the form $[Y, Y^{3/2}]$, and the interface
demands that such partial information suffice to exclude zeros.

The formal, precise statement, with the complete list of validity
conditions on the certificate and the complete formulas for all
thresholds, is Theorem~\ref{thm:main} in \S\ref{sec:main}. The algorithm
that computes the address is given in \S\ref{sec:algorithm} and specified
for implementers in the companion file \texttt{certificates.md}. A
self-contained sanity check of the constants is planned for
\texttt{numerics/}.

\subsection{Design decisions, made explicit}\label{sec:design}

\textbf{Hypothesis injection.} We never produce zeros of $\zeta$, and we
never verify zero-free regions ourselves. The certificate \emph{asserts}
(upstream, by interval arithmetic, in the style of rigorous zero
computations) that a certain slab contains exactly one zero of $\zeta$ to
the right of the critical line, inside a certain rectangle. Everything we
do is conditional on that assertion; the assertion is a hypothesis, never
an axiom. This follows the discipline of recent explicit number theory in
proof assistants, and it is what keeps the eventual PA-track (see
\S\ref{sec:pa}) low: the transcendental content enters as a hypothesis of
arithmetic form.

\textbf{One zero at a time.} The certificate isolates a \emph{single}
zero. Distinct off-line zeros that lie within the (necessarily tall,
see~\S\ref{sec:limits}) certified zero-free slab of one another cannot be
handled by the present argument, which needs the isolated pole's
principal part to dominate uncontested; see \S\ref{sec:limits} for what is
lost and how it might be recovered. A multiple zero at a single point is
not an obstruction (the pole of the Mellin transform is then simple with
residue scaled by multiplicity) but our certificate as formulated demands
a simple zero; the generalization is routine and omitted.

\textbf{Terrible constants, on purpose.} We track constants with
deliberately safe margins (e.g.\ $\Sig = 252$ where any sufficiently large
integer works, $4x^2$ where $x^2$ is true for large $x$). The point of
this paper is computability, not competitiveness; a
verification-grade version is a natural follow-up with the community that
does that for a living. Every displayed constant is a rational number.

\textbf{Quoted standard estimates.} Three inputs are quoted from the
classical literature rather than re-derived with our own constants
(Lemmas~\ref{lem:zeta-prime-prime}, \ref{lem:zlow}, \ref{lem:backlund}):
an explicit partial-fraction bound for $\zeta'/\zeta$ away from zeros; the
classical low-lying zero data (first zero at height $14.1347\ldots$, no
zeros below, first zeros on the line); and an explicit local zero count.
We state them with conservative constants, flag them as such, and note
that any explicit version suffices for every use we make of them. The
effective-shadow work (Phase 2) will need to re-prove them anyway, and
that is the natural home for tightening.

\textbf{Honesty about the Alp\"oge claim.} The announced Lean-formalized
proof of RH \cite{Alp26} was independently machine-verified on November
21, 2026 \cite{Rub26v}: a public end-to-end compile of all $175{,}187$
modules, an axiom audit reducing the five headline theorems to Lean's
three standard axioms, a full replay of every module through the
kernel, and a check of all $1{,}484{,}729$ relevant declarations by a
second, independently written kernel; independent confirmations by
Sutherland (MIT), Arledge, and DeDeo are tracked on Naskr\k{e}cki's
\texttt{riemann-rev} dashboard. This settles the mechanical correctness
of the development from Lean's axioms; the new mathematics inside it
will take the community a long time to digest at human scale, and
nothing in this paper depends on any of it. The windowed lemma
(Lemma~\ref{lem:rubinstein}) and its one-page proof exist in the
classical exposition \cite{Rub26} (independently of the archive),
credited there; our contribution is the explicitness, the certificate
formulation, and the computable address.

\subsection{Overview of the argument}

Write $u = \log x$, $g(u) = G(e^{u})$, and let
\begin{equation}\label{eq:Fdef}
\begin{split}
F(s) \;&=\; \int_1^\infty G(x)\, x^{-s-1}\,dx
      \;=\; \int_0^\infty g(u)\, e^{-su}\,du\\
     \;&=\; -\frac{\zeta'(s)}{\zeta(s)} \cdot \frac{1}{s(s-1)}
            \;-\; \frac{1}{(s-1)^2} \;+\; \frac{1+\Euler}{s-1},
\end{split}
\end{equation}
the Mellin transform of $G$, meromorphic on $\C$ with simple poles at the
nontrivial zeros $\rho$ of $\zeta$ (residue $-\nu_{\rho}/(\rho(\rho-1))$,
with $\nu_\rho$ the multiplicity), at $s=0$, and at the trivial zeros.
Let
\[
A_j(s) \;=\; \frac{(-1)^j F^{(j)}(s)}{j!},
\qquad
K_j(u) \;=\; \frac{u^j e^{-\Sig u}}{j!},
\]
so that for $\re s = \Sig > 2$ and real $t$,
$A_j(\Sig + it) = \int_0^\infty g(u) e^{-itu} K_j(u)\,du$.
The kernel $K_j$ peaks at $u = j/\Sig$; with $U = \log Y$ and $j$ the
nearest integer to $1.23(\Sig-2)U$, the peak sits in the middle of the
window $[U, \tfrac32 U]$ and the mass of $K_j$ outside the window is
exponentially small (Lemma~\ref{lem:tail}), with decay $e^{-cj}$ paying
exactly for the conversion between the scales $(\Sig-2)$ and
$(\Sig-\theta)$ at the cost constant $e^{c}$ --- this is the numerical
content of Rubinstein's choices $\eta$, $c$, $\Sig \ge 252$.

The ceiling hypothesis $G \le C'x^{\theta}$ on the window, together with
the crude unconditional bound $|G(x)| \le 4x^2$ (\S\ref{sec:statistic}),
controls the \emph{positive part} of $g$, hence (via
$|g| + g = 2\max(g,0)$, applied to the pair of points
$\Sig+it$, $\Sig$) controls $|A_j(\Sig+it)| + A_j(\Sig)$.
A Cauchy estimate at the real point controls $A_j(\Sig)$
(Lemma~\ref{lem:sigcauchy}), so $|A_j(z)|$ itself is small for every real
$t$: small like $(\Sig-\theta)^{-j}$.

On the other hand, the certified zero $\rho$ with $\beta = \re\rho > \theta$
is a pole of $F$ at distance $d \approx \Sig - \beta < \Sig - \theta$ from
the explicit point $z = \Sig + i\,\mathrm{mid}(R)$ on the line
$\re s = \Sig$; subtracting its principal part and applying Cauchy on a
circle that the certificate guarantees contains no other pole
(Lemma~\ref{lem:pole-lower}) gives $|A_j(z)| \ge c_1 d^{-j-1}$, which
beats $(\Sig-\theta)^{-j}$ by the exponential factor
$((\Sig-\theta)/d)^{j+1}$.

The contradiction is made \emph{quantitative}: if the ceiling holds with
constant $C'$, then $C' \ge Q(j)$ where $Q(j)$ is an explicit huge number
growing like $((\Sig-\theta)/d)^{j}$ (Theorem~\ref{thm:main}); since $Q(j)
\to \infty$ along $j \sim 1.23(\Sig-2)\log Y$, the ceiling must fail, and
by continuity of $G$ it must fail at an attainable maximum of $G(x)/x^\theta$
over the window, with margin $m = Q(j) - C$. A grid search with explicit
mesh, evaluating $G$ by rigorous interval arithmetic, then \emph{finds}
the address (Proposition~\ref{prop:search}), with a proof of termination;
the total work is polynomial in $Y$.

Everything needed for this is contained in \S\ref{sec:statistic}--\S
\ref{sec:main}. \S\ref{sec:algorithm} specifies the certificate checker's
algorithm. \S\ref{sec:limits} states limitations (clusters, quoted
inputs), \S\ref{sec:lfunction} sketches the Dirichlet $L$-function
variant, \S\ref{sec:pa} records notes toward the Lean two-layer
formalization and the PA track, and \S\ref{sec:registry} is the constant
registry. Appendix~\ref{sec:app:elementary} collects the elementary
inequalities with proofs.

\subsection{Acknowledgments}

The windowed lemma and its one-page proof are due to Michael Rubinstein's
explanatory draft \cite{Rub26} (with AI co-authors), which in turn
explains the Alp\"oge development \cite{Alp26}; the classical skeleton is
Landau's \cite{Lan05} with the P\'olya--Ingham kernel method
\cite{Ing32}. We claim credit only for the effectivization, the
certificate formulation, and the computational track. The layman-facing
exposition we mirror is the segment ``Step two: the ceiling'' of
\cite{RieVideo}.

\section{The statistic and its Mellin transform}\label{sec:statistic}

\subsection{Notation}

$\Lam(n)$ is the von Mangoldt function: $\log p$ if $n = p^k$ is a prime
power, $0$ otherwise. $\psi(x) = \sum_{n\le x}\Lam(n)$. $\Euler$ is the
Euler--Mascheroni constant. $\zeta$ is the Riemann zeta function;
``zeros'' always means nontrivial zeros, which satisfy
$0 < \beta < 1$ and (since $\zeta(\sigma)<0$ on $0<\sigma<1$) are
non-real. The classical first zero has height
$\gamma_1 = 14.134725\ldots$, there are no zeros with
$0<|\im s| < \gamma_1$, and the first zeros lie on the critical line
(Lemma~\ref{lem:zlow}). $e$ and $\log$ are the natural exponential and
logarithm; $x^{\theta}$ for real $x>0$ is $e^{\theta\log x}$.

Throughout the paper $\Sig = 252$, $c = 0.008$ and
\begin{equation}\label{eq:sigma-c}
\frac{\Sig - \theta}{\Sig - 2} \;\le\; e^{c}
\qquad\text{for every } 0 < \theta < 1,
\end{equation}
since $\Sig - \theta \le 252 < 250\,e^{0.008} = 252.008\ldots \le
(\Sig-2)e^{c}$.

\subsection{Elementary size of the statistic}

\begin{lemma}[$G$ is small and piecewise linear]\label{lem:G-size}\
\begin{enumerate}
\item For all $x \ge 1$: $|G(x)| \le 4x^2$.
\item $G$ is continuous on $[1,\infty)$; on the interval between two
consecutive prime powers it is linear in $x$; in particular $G$ is
Lipschitz on any bounded interval with constant
$L_G(x) := \tfrac12(\log x)^2 + 2\log x + 1.6$ on $[2, x]$:
$|G(x_2)-G(x_1)| \le L_G(x_2)\,|x_2-x_1|$ for $2 \le x_1 \le x_2$.
\end{enumerate}
\end{lemma}

\begin{proof}
(1) First, $S(x) = x\int_1^x \psi(t)\,t^{-2}\,dt$: each term
$\Lam(n)(x/n - 1)$ equals $\Lam(n)\,x\int_n^x t^{-2}\,dt$, and the sum
over $n \le t$ is $\psi(t)$. Since $\psi(t) \le \sum_{n \le t}\log n \le
t\log t$ (and $\psi(t) = 0$ for $t < 2$),
\[
S(x) \;\le\; x \int_2^x \frac{\log t}{t}\,dt
\;=\; \frac{x}{2}\bigl((\log x)^2 - (\log 2)^2\bigr) \;\le\;
\frac{x(\log x)^2}{2}.
\]
Also $x\log x \le x(\log x)^2$ and $(1+\Euler)x \le 1.6\,x(\log x)^2$
for $x \ge e$. Hence $|G(x)| \le 3.1\,x(\log x)^2$ for $x \ge e$, and
$3.1(\log x)^2 \le 4x$ there (the function $4x - 3.1(\log x)^2$
increases for $x \ge e$ and is positive at $x=e$). For $1 \le x < e$ we
have $S(x) \le \Lam(2)(x/2 - 1) \le 0.35$ (only $n=2$ can contribute),
so $|G(x)| \le 0.35 + x\log x + (1+\Euler)x \le 0.35 + e + 1.6e \le
7.1 \le 4x^2$ on $[1, 2]$, and on $[2,e]$, $|G(x)| \le 0.35 + e + 1.6e
\le 7.1 \le 16 \le 4x^2$. This proves (1).

(2) $S$ is continuous: as $x$ crosses a prime power $n$, the summand
$\Lam(n)(x/n-1)$ switches on and vanishes at $x = n$. Between prime
powers, $S(x) = x A(x) - \psi(x)$ with $A(x) = \sum_{n\le x}\Lam(n)/n$
and $\psi(x)$ constant, so $S$ is linear there and so is $G$, with slope
$G'(x) = A(x) - \log x + 1 + \Euler$. Finally
\[
A(x) \;=\; \frac{\psi(x)}{x} + \int_1^x \frac{\psi(t)}{t^2}\,dt
\;\le\; \log x + \frac{(\log x)^2}{2},
\]
so $|G'| \le \tfrac12(\log x)^2 + 2\log x + 1.6$ where defined, and the
Lipschitz claim follows since a continuous, piecewise linear function is
Lipschitz with any upper bound for its slopes.
\end{proof}

\subsection{The Mellin transform and its poles}

\begin{proposition}\label{prop:mellin}\
\begin{enumerate}
\item For $\re s > 1$ the integral \eqref{eq:Fdef} converges and
\[
\begin{split}
F(s) \;&=\; \sum_{n\ge1} \frac{\Lam(n)\,n^{-s}}{s(s-1)}
\;-\; \frac{1}{(s-1)^2} \;+\; \frac{1+\Euler}{s-1}\\
     \;&=\; -\frac{\zeta'(s)}{\zeta(s)}\cdot\frac{1}{s(s-1)}
\;-\; \frac{1}{(s-1)^2} \;+\; \frac{1+\Euler}{s-1}.
\end{split}
\]
\item The right-hand side is meromorphic on $\C$; at a zero $\rho$ of
$\zeta$ of multiplicity $\nu$, $F$ has a simple pole with residue
\[
\res_{s=\rho} F \;=\; -\frac{\nu}{\rho(\rho-1)}.
\]
$F$ is holomorphic at $s = 1$ (the principal parts cancel), and has poles
at $s=0$ and at the trivial zeros $-2, -4, \ldots$
\end{enumerate}
\end{proposition}

\begin{proof}
For $\re s > 1$, termwise integration gives
$\int_1^\infty (x/n-1)^+ x^{-s-1}dx = n^{-s}/(s(s-1))$ (substitute
$x = nt$), so the transform of $S$ is
$\frac{1}{s(s-1)}\sum \Lam(n)n^{-s} = -\zeta'(s)/\zeta(s) \cdot
\frac{1}{s(s-1)}$; the transform of $-x\log x$ is $-\int_1^\infty
x^{-s}\log x \,dx = -1/(s-1)^2$ and that of $(1+\Euler)x$ is
$(1+\Euler)/(s-1)$. Convergence of $\int_1^\infty G(x)x^{-s-1}dx$ for
$\re s > 1$ follows from Lemma~\ref{lem:G-size} and
$S(x) \le \tfrac12 x(\log x)^2$.

(2) $\zeta'/\zeta$ is meromorphic with simple poles of residue $\nu$ at
each zero $\rho$; multiplying by the holomorphic non-vanishing factor
$1/(s(s-1))$ (non-vanishing because $\rho \notin \{0,1\}$) gives a simple
pole of residue $-\nu/(\rho(\rho-1))$ at $\rho$ and a pole at $s=0$.
At $s = 1$: $\zeta(s) = \frac{1}{s-1} + \Euler + O(s-1)$, so
$-\zeta'/\zeta \cdot \frac{1}{s(s-1)} = \frac{1}{(s-1)^2} -
\frac{1+\Euler}{s-1} + O(1)$, which cancels the two remaining terms.
\end{proof}

\begin{remark}
We never use the explicit formula for $G$; only the transform identity
\eqref{eq:Fdef}, its meromorphic continuation, and the residue formula
enter the proof.
\end{remark}

\begin{lemma}[differentiation under the integral]\label{lem:diff}
For $\re s \ge \Sig$ (indeed $\re s > 2$) and $j \ge 0$,
\[
A_j(s) \;=\; \int_0^\infty (-u)^j g(u) e^{-su}\,\frac{du}{j!},
\]
the integral being absolutely convergent. In particular, for real $t$,
\begin{equation}\label{eq:Atau}
A_j(\Sig + it) \;=\; \int_0^\infty g(u)\, e^{-itu}\, K_j(u)\,du.
\end{equation}
\end{lemma}

\begin{proof}
By Lemma~\ref{lem:G-size}(1), $|g(u)| \le 4e^{2u}$, so
$u^j g(u) e^{-su}$ is dominated by $4u^j e^{-(\re s - 2)u}$, integrable
for $\re s > 2$; differentiate $j$ times under the integral sign.
\end{proof}

\section{The kernel and its window}\label{sec:kernel}

Fix for the rest of the paper:
\[
\Sig = 252, \qquad c = 0.008, \qquad \eta_* = 0.983, \qquad
[v_-, v_+] = [0.88,\, 1.12],
\]
and recall \eqref{eq:sigma-c}. All four are rationals; the roles are:
$\Sig$ the abscissa at which derivatives are taken, $c$ the tail decay
rate (which also pays for the scale conversion), $\eta_*$ the cap on the
P\'olyya kernel outside the window, and $[v_-,v_+]$ the fixed interval
about the peak inside which the kernel is large.

\begin{lemma}[choice of $j$]\label{lem:j-range}
Let $U \ge 3$ and let $j$ be the nearest integer to $1.23(\Sig-2)U$.
Then $j \ge 1$ and, writing
\[
v_1 = \frac{(\Sig-2)U}{j}, \qquad v_2 = \frac{\tfrac32(\Sig-2)U}{j},
\]
we have
\begin{equation}\label{eq:vwindow}
0 < v_1 \;\le\; 0.825, \qquad v_2 \;\ge\; 1.2,
\end{equation}
and $[v_-,v_+] \subseteq [v_1, v_2]$.
\end{lemma}

\begin{proof}
$j \ge 1.23 \cdot 250 \cdot 3 - \tfrac12 > 921$. For the upper bound,
$j \le 1.23(\Sig-2)U + \tfrac12 \le 1.25(\Sig-2)U$ iff $0.02 \cdot 250
\cdot U \ge \tfrac12$, i.e.\ $U \ge 0.1$; then $v_1 \ge 0.8$. For the
lower bound, $j \ge 1.23(\Sig-2)U - \tfrac12 \ge 1.2122(\Sig-2)U$ iff
$0.0178 \cdot 250 U \ge \tfrac12$, i.e.\ $U \ge 0.112$; then $v_1 \le
1/1.2122 = 0.8249 \le 0.825$ and $v_2 = 1.5/1.2122$-to-$1.25 \ge 1.2$.
Since $v_1 \le 0.825 < 0.88$ and $v_2 \ge 1.2 > 1.12$, the last claim
holds.
\end{proof}

\begin{lemma}[tail of the P\'olya kernel]\label{lem:tail}
Let $U \ge 3$, $j$ as in Lemma~\ref{lem:j-range}, and $v_1, v_2$ as in
\eqref{eq:vwindow}. Put $b(v) = v e^{1-v}$ (so $b(1)=1$, $b$ increasing
on $[0,1]$, decreasing on $[1,\infty)$). Then:
\begin{enumerate}
\item (\emph{outside is capped}) $b(v_1) \le b(0.825) < 0.983 = \eta_*$
and $b(v_2) \le b(1.2) < 0.983 = \eta_*$; and for $v \notin [v_1,v_2]$
(within $[0,\infty)$) we have $b(v) \le \eta_*$.
\item (\emph{inside is large}) $b(v) \ge 0.9921$ for all
$v \in [v_-,v_+] = [0.88,1.12]$.
\item (\emph{fraction bound})
\[
\frac{\displaystyle\int_{[0,v_1]\cup[v_2,\infty)} b(v)^j\,dv}
     {\displaystyle\int_0^\infty b(v)^j\,dv}
\;\le\; 13\, e^{-cj}.
\]
\end{enumerate}
\end{lemma}

\begin{proof}
(1) Numerically (all values rational-rounded to safety; see
Appendix~\ref{sec:app:elementary} for the displayed arithmetic):
$b(0.825) = 0.825 e^{0.175} \le 0.983$, $b(1.2) = 1.2e^{-0.2} \le
0.983$. On $[0,v_1]$, $b$ is increasing, so $b \le b(v_1) \le \eta_*$; on
$[v_2,\infty)$, $b$ is decreasing, so $b \le b(v_2) \le \eta_*$.

(2) $b(0.88) = 0.88e^{0.12} \ge 0.9921$ and $b(1.12) = 1.12e^{-0.12}
\ge 0.9921$, and $b \ge \min\{b(0.88), b(1.12)\}$ on $[0.88,1.12]$
because $b$ increases then decreases with its peak at $1$ inside the
interval.

(3) Numerator: $\int_0^{v_1} b^j \le v_1 b(v_1)^j \le
0.83\,\eta_*^j$
(monotonicity); $\int_{v_2}^{2v_2} b^j \le v_2 \eta_*^j \le
1.24\,\eta_*^j$ (here $v_2 \le \tfrac{1.5}{1.2122} < 1.24$ by the proof
of Lemma~\ref{lem:j-range}); and for $v \ge 2v_2 \ge 2.4$, writing
$w = v/2 \ge 1.2$ we have $b(v) = 2\,b(w)\,e^{-w} < 2\,e^{-w} =
2\,e^{-v/2}$ (since $b(w) \le b(1.2) < 1$), hence
$b(v)^j < 2^j e^{-jv/2}$ and, shifting the origin to $2.4$,
$b(v)^j < (2e^{-1.2})^j e^{-j(v-2.4)/2} < (0.61)^j e^{-j(v-2.4)/2}$,
so $\int_{2v_2}^\infty b^j\,dv \le 0.61^j
\int_0^\infty e^{-jv/2}\,dv = 0.61^j\,\frac{2}{j} \le \eta_*^j$ for
$j \ge 4$. Total numerator $\le (0.83 + 1.24 + 1)\,\eta_*^j \le
3.1\,\eta_*^j$.
Denominator $\ge \int_{v_-}^{v_+} b^j \ge 0.24 \cdot (0.9921)^j$.
Hence the fraction is at most
$\frac{3.1}{0.24}\bigl(\frac{0.983}{0.9921}\bigr)^{j} < 12.92 \cdot
(0.99083)^{j} \le 13\,e^{-0.008j}$, since $0.99083 \le e^{-0.008} =
0.99203$. This holds for every $j \ge 4$, and our $j$ exceeds $900$.
\end{proof}

\begin{remark}
The constant pair $(\eta_*, c)$ mirrors Rubinstein's choices
\begin{gather*}
\eta = \max\bigl(0.8e^{0.2},\ 1.2e^{-0.2}\bigr) = 0.9825\ldots,\\
c = 0.008 < \log\frac{1+\eta}{2\eta} = 0.00887\ldots,
\end{gather*}
in that the crucial inequality for our constants is
$\frac{2\eta_*}{1+\eta_*} \le e^{-c}$, i.e.\ the capped outside loses
to the inside at exponential rate at least $c$. We chose slightly
worse $\eta_* = 0.983$ than $\eta$ to buy rounding room.
\end{remark}

\begin{remark}[the window exponent is symbolic]\label{rem:exponent}
The architecture's window exponent $3/2$ enters only through the
numbers in this section. For a general window $[Y, Y^{\omega}]$
with $1 < \omega < 2$, the kernel range is $j \in
[(\Sig-2)U/v_L^{\omega}, (\Sig-2)U\omega/v_R^{\omega}]$ with the
$v$-window $[v_L^{\omega}, v_R^{\omega}] := [1/\sqrt{\omega},
\sqrt{\omega}]$ (plus rounding), the outside cap becomes
\[
\eta_{\omega} := \max\bigl(b(1/\sqrt{\omega}),\; b(\sqrt{\omega})\bigr)
= \max_{\epsilon > 0} \text{rounding margins},
\]
any $c_{\omega} < \log\frac{1+\eta_{\omega}}{2\eta_{\omega}}$ replaces
$c$, and $\Sig$ is re-chosen as the least integer with
$(\Sig-\theta)/(\Sig-2) \le e^{c_{\omega}}$ for all $\theta \in (0,1)$.
For $\omega = 3/2$: $1/\sqrt{\omega} \approx 0.816$,
$\sqrt{\omega} \approx 1.225$, $\eta_{\omega} \approx 0.981$, and the
choices above are compatible with these values plus rounding room. The
rest of the paper is unchanged; we fix $\omega = 3/2$ because that is
the architecture's constant (the $L$-function side of the archive
additionally uses $\omega = 2$, for which the same formulas apply).
\end{remark}

\begin{lemma}[mass outside the window]\label{lem:tail2}
Let $U \ge 3$, $j$ as in Lemma~\ref{lem:j-range}. Then, for any
(measurable) set $E \subseteq [0,\infty) \setminus [U,
\tfrac32 U]$,
\[
\int_E 4\,e^{2u} K_j(u)\,du \;\le\; 52\, e^{-cj}\,(\Sig-2)^{-j-1}
\;\le\; 52\, e^{c}\,(\Sig-\theta)^{-j-1},
\]
the second inequality by \eqref{eq:sigma-c}.
\end{lemma}

\begin{proof}
Substitute $u = jv/(\Sig-2)$, $du = j\,dv/(\Sig-2)$: then
$u^j e^{-(\Sig-2)u}\,du = \bigl(\tfrac{j}{\Sig-2}\bigr)^{j+1} e^{-j}
b(v)^j\,dv$, the window $[U, \tfrac32 U]$ becomes $[v_1, v_2]$, and the
normalizing constant is the same for the full integral, which equals
$\int_0^\infty u^j e^{-(\Sig-2)u}du = j!\,(\Sig-2)^{-j-1}$. Hence
\[
\begin{split}
\int_E 4\,e^{2u}K_j(u)\,du
\;&=\; 4\,(\Sig-2)^{-j-1}\cdot
  \frac{\int_{E'} b(v)^j\,dv}{\int_0^\infty b(v)^j\,dv}\\
\;&\le\; 4 \cdot 13\,e^{-cj}\,(\Sig-2)^{-j-1}
\;=\; 52\,e^{-cj}(\Sig-2)^{-j-1},
\end{split}
\]
where $E'$ is the image of $E$ under the substitution and we used
Lemma~\ref{lem:tail}(3).
\end{proof}

\begin{lemma}[mass inside the window]\label{lem:inside}
For any $\theta' \ge 0$ and any $E \subseteq [0,\infty)$,
\[
\int_E e^{\theta' u} K_j(u)\,du \;\le\; (\Sig-\theta')^{-j-1},
\]
whenever $\theta' < \Sig$.
\end{lemma}

\begin{proof}
$\int_E e^{\theta'u}u^j e^{-\Sig u}\frac{du}{j!} \le \int_0^\infty
u^j e^{-(\Sig-\theta')u}\frac{du}{j!} = (\Sig-\theta')^{-j-1}$.
\end{proof}

\section{The one-sided Mellin bound}\label{sec:onesided}

Throughout this section, fix $U \ge 3$ (so $Y = e^U \ge 20$), let $j$ be
the nearest integer to $1.23(\Sig-2)U$, and let $C' \ge 0$ be arbitrary.
Recall $K_j(u) = u^j e^{-\Sig u}/j!$ and \eqref{eq:Atau}.

\begin{lemma}[one-sided control]\label{lem:onesided}
Suppose
\[
G(x) \;\le\; C' x^{\theta} \qquad \text{for all } x \in [Y, Y^{3/2}].
\]
Then for every real $t$,
\begin{equation}\label{eq:onesided}
|A_j(\Sig + it)| \;+\; A_j(\Sig)
\;\le\; \bigl(2C' + C_1\bigr)\,(\Sig-\theta)^{-j-1},
\qquad C_1 := 105.
\end{equation}
\end{lemma}

\begin{proof}
By \eqref{eq:Atau}, for real $t$,
$|A_j(\Sig+it)| \le \int_0^\infty |g(u)|K_j(u)\,du$ and
$A_j(\Sig) = \int_0^\infty g(u) K_j(u)\,du$ (both integrands real).
Using $|g| + g = 2\max(g, 0) = 2g^+$ pointwise,
\[
|A_j(\Sig+it)| + A_j(\Sig)
\;\le\; \int_0^\infty (|g(u)| + g(u))K_j(u)\,du
\;=\; 2\int_0^\infty g^+(u)K_j(u)\,du.
\]
Split the last integral at the window $[U, \tfrac32 U]$. On the window,
$g^+(u) = \max(G(e^u),0) \le C' e^{\theta u}$ by hypothesis, so
Lemma~\ref{lem:inside} bounds that part by $C'(\Sig-\theta)^{-j-1}$.
Off the window, $g^+ \le |g| \le 4e^{2u}$ by
Lemma~\ref{lem:G-size}(1) (recall $u = \log x$), so
Lemma~\ref{lem:tail2} bounds that part by
$52 e^{c}(\Sig-\theta)^{-j-1} \le 52.5\,(\Sig-\theta)^{-j-1}$.
Total: $|A_j(\Sig+it)| + A_j(\Sig) \le (2C' + 105)(\Sig-\theta)^{-j-1}$.
\end{proof}

\begin{lemma}[Cauchy at the real point]\label{lem:sigcauchy}
Let $r' \in (0, \Sig)$ be such that $F$ is holomorphic on an open set
containing the closed disk $D(\Sig, r') = \{w : |w - \Sig| \le r'\}$. Put
$M' := \sup_{|w-\Sig| = r'}|F(w)|$. Then for all $j \ge 0$,
\[
|A_j(\Sig)| \;\le\; M' \, (r')^{-j}.
\]
\end{lemma}

\begin{proof}
$A_j(\Sig) = (-1)^j [w^j]F(\Sig + w)$ and
$|[w^j]F(\Sig+w)| \le M' (r')^{-j}$ by Cauchy's coefficient estimate.
\end{proof}

\begin{corollary}\label{cor:upper}
Under the hypotheses of Lemma~\ref{lem:onesided}, if moreover $r'$ and
$M'$ are as in Lemma~\ref{lem:sigcauchy}, then for every real $t$,
\begin{equation}\label{eq:upper}
|A_j(\Sig + it)| \;\le\; (2C' + C_1)(\Sig-\theta)^{-j-1} + M'(r')^{-j}.
\end{equation}
\end{corollary}

\begin{proof}
$|A_j(\Sig+it)| \le (2C'+C_1)(\Sig-\theta)^{-j-1} - A_j(\Sig) \le
(2C'+C_1)(\Sig-\theta)^{-j-1} + |A_j(\Sig)|$.
\end{proof}

\section{The certificate and the pole}\label{sec:certificate}

\subsection{Certificate data and validity conditions}

\begin{definition}[ceiling-refutation certificate]\label{def:cert}
A \emph{certificate} is a tuple of rational numbers
\[
(\theta,\; C,\; \beta_1, \beta_2, \gamma_1, \gamma_2,\; h)
\]
with
\[
\theta \in [\tfrac12, 1), \qquad C \ge 1, \qquad
\beta_1 < \beta_2, \qquad \gamma_1 < \gamma_2, \qquad h > 0,
\]
subject to the \emph{validity conditions}
\begin{align}
\text{(V1)}&\quad \theta < \beta_1 \le \beta_2 < 1, \notag\\[2pt]
\text{(V2)}&\quad (\Sig - \beta_1)^2 + \frac{\Gamma^2}{4}
        < (\Sig - \theta)^2,
        \qquad \Gamma := \gamma_2 - \gamma_1, \notag\\[2pt]
\text{(V3)}&\quad (\Sig - \beta_1)^2 + \frac{\Gamma^2}{4}
        < (\Sig - \beta_2)^2 + \gamma_1^2, \notag\\[2pt]
\text{(V4)}&\quad h \;\ge\; \frac{\Gamma}{2} + 17, \notag\\[2pt]
\text{(V5)}&\quad \min\bigl(\varepsilon_z,\ \delta',\ W_1,\ W_1'\bigr)
\;\ge\; 10^{-6}, \notag
\end{align}
Here $\varepsilon_z, \delta', W_1, W_1'$ are as defined in
\S\ref{sec:derived}, so (V5) is a finite arithmetic check.
together with the \emph{certified statement}, produced upstream (not by
us) by rigorous computation:
\begin{equation}\tag{S1}\label{eq:S1}
\begin{minipage}{0.86\linewidth}
The open slab
$\Sigma := \{s \in \C : \re s > \tfrac12,\;
\gamma_1 - h < \im s < \gamma_2 + h\}$
contains exactly one zero $\rho$ of $\zeta$, counted with multiplicity;
that zero is simple and lies in the open rectangle
$R := [\beta_1, \beta_2] \times [\gamma_1, \gamma_2]$.
\end{minipage}
\end{equation}
\end{definition}

\begin{remark}[what each condition buys]\label{rem:conditions}\
\begin{itemize}
\item (V1) says the alleged zero lies strictly to the right of the
ceiling exponent; it makes the amplification factor
$\lambda = (\Sig-\theta)/\dbar$ (below) strictly exceed $1$.
\item (V2) says the rectangle is short enough (or the zero far enough
right) that the \emph{entire rectangle}, seen from the point
$z = \Sig + i\,\mathrm{mid}$ (below), is strictly closer to $z$ than the
critical line is; note (V2) implies
$\dbar < \Sig - \tfrac12 = 251.5$, since $\theta \ge \tfrac12$.
\item (V3) is the same comparison against the \emph{nearest point of the
rectangle to the real point} $\Sig$: it guarantees that the Cauchy
circle used at $\Sig$ (radius $r'$, below) can be placed strictly between
$\dbar$ and the distance from $\Sig$ to the certified zero.
\item (V4) makes the slab tall enough to contain everything the second
Cauchy circle ``sees'' of the critical strip (reach $\le \sqrt{251.5^2
- 251^2} < 15.9$ around the height of $z$), with one unit of margin for
the local zero count.
\item (S1) is the hypothesis-injected content: it converts the classical
``choose $t$ so the nearest pole is unique'' into an explicit, checkable
input. It is exactly the kind of statement that rigorous zero
computations produce (argument principle plus interval arithmetic).
\end{itemize}
Conditions (V2) and (V3) restrict how thick the rectangle may be,
roughly $\Gamma < 2\sqrt{(\beta_1 - \tfrac12)\cdot 503}$ and
$\beta_2 - \beta_1 < \gamma_1^2/503$; both are mild for realistic
isolation rectangles, and both are decidable by trivial rational
arithmetic.
\end{remark}

\subsection{Derived quantities}\label{sec:derived}

With the certificate fixed, define the following rational-computable
quantities (square roots computed to any needed precision by rational
interval arithmetic; the theorem only ever uses them through explicit
rational upper/lower bounds, so any certified approximation suffices):
\begin{align}
\Gamma &:= \gamma_2 - \gamma_1, &
\mathrm{mid} &:= \tfrac12(\gamma_1 + \gamma_2), \notag\\
z &:= \Sig + i\,\mathrm{mid}, &
\dbar &:= \sqrt{(\Sig - \beta_1)^2 + \tfrac{\Gamma^2}{4}}, \notag\\
M_1 &:= \sqrt{\beta_2^2 + \gamma_2^2}, &
M_2 &:= \sqrt{(1-\beta_1)^2 + \gamma_2^2}, \notag\\
D_R &:= \sqrt{(\Sig - \beta_2)^2 + \gamma_1^2}, &
\gamma_{\mathrm{first}} &:= 14.1347, \notag\\
D_{\mathrm{line}} &:= \sqrt{(\Sig - \tfrac12)^2 +
\gamma_{\mathrm{first}}^2} > 251.89, &
D_{\min} &:= \min(D_{\mathrm{line}},\, D_R), \notag\\
\lambda &:= \frac{\Sig - \theta}{\dbar} > 1, &
r_z &:= \tfrac12\bigl(\dbar + (\Sig - \tfrac12)\bigr), \notag\\
\varepsilon_z &:= (\Sig-\tfrac12) - r_z > 0, \notag\\
r' &:= \tfrac12(\dbar + D_{\min}), &
\delta' &:= D_{\min} - r' > 0, \notag
\end{align}
so that $\varepsilon_z = \tfrac12\bigl((\Sig-\tfrac12) - \dbar\bigr)$
and $\delta' = \tfrac12(D_{\min} - \dbar)$.
Note that $\dbar = \max_{w \in R}|z - w|$: the farthest corner of $R$
from $z$ is $(\beta_1, \gamma_1)$ or $(\beta_1, \gamma_2)$, and its
distance is $\sqrt{(\Sig-\beta_1)^2 + \Gamma^2/4}$. Also $M_1 M_2 \ge
|\rho(\rho-1)|$ for every $\rho \in R$, so that by
Proposition~\ref{prop:mellin}(2),
\begin{equation}\label{eq:residue}
|c_{\rho}| \;=\; \frac{1}{|\rho(\rho-1)|} \;\ge\; \frac{1}{M_1 M_2},
\qquad c_\rho := \res_{s=\rho}F,
\end{equation}
since $\rho$ is simple. The validity conditions give
\[
\text{(V2): } \dbar < \Sig-\tfrac12; \qquad
\text{(V3) and } \dbar < 251.5 < D_{\mathrm{line}}:\quad
\dbar < D_{\min}.
\]
Hence $r_z < \Sig - \tfrac12 = 251.5$, $\varepsilon_z > 0$, $r' > \dbar$,
$\delta' > 0$, and $\lambda > 1$. Finally put
\begin{align*}
W_0 &:= \tfrac12, &
W_1 &:= \bigl||z-1| - r_z\bigr|, &
W_0' &:= 252 - r', \\
W_1' &:= r' - 251, &
t_{\max} &:= \max(|\gamma_1|, |\gamma_2|) + 252,
\end{align*}
so that $r_z < 251.5 < 252$ gives, for every $w$ on the $z$-circle,
\[
|\im w| \;\le\; |\mathrm{mid}| + r_z \;<\; \max(|\gamma_1|,|\gamma_2|) + 252
\;=\; t_{\max},
\]
and (V5) asserts $\min(\varepsilon_z, \delta', W_1, W_1') \ge 10^{-6}$.

\subsection{The two Cauchy circles are pole-free}

We record the classical low-zero inputs once.

\begin{assumption}[low zero data; classical]\label{lem:zlow}\
\begin{enumerate}
\item $\zeta(\sigma) < 0$ for $0 < \sigma < 1$; in particular $\zeta$ has
no real zeros in the critical strip, and every zero is non-real.
\item $\zeta$ has no zeros with $0 < |\im s| < 14.1347$; the first zero
lies at $\tfrac12 \pm 14.134725\ldots i$.
\item Every zero with $0 < |\im s| \le 30$ has $\re s = \tfrac12$.
\end{enumerate}
Items (2), (3) are classical rigorous zero computations: Lehman
\cite{Leh66} located the first zeros on the line; Turing \cite{Tur53}
and Brent \cite{Bre79} verified the first zeros; Platt--Trudgian
\cite{PT21} verified the Riemann hypothesis to height $3\cdot 10^{12}$,
which subsumes (2) and (3) at our scales with enormous margin. Item (1)
is elementary (alternating series).
\end{assumption}

\begin{lemma}[the circles avoid all poles of $F$]\label{lem:circles}\
\begin{enumerate}
\item Every zero $\rho'$ of $\zeta$ with $\rho' \ne \rho$ satisfies
\[
|z - \rho'| \;\ge\; \Sig - \tfrac12 \;=\; 251.5 \;>\; r_z.
\]
Consequently the closed disk $D(z, r_z)$ contains no zero of $\zeta$
other than $\rho$, and $\rho$ itself satisfies
$|z - \rho| \le \dbar < r_z$; moreover every point $w$ with
$|w - z| = r_z$ has $\re w > \Sig - r_z > \tfrac12$ and satisfies
\[
|w - \rho'| \;\ge\; \varepsilon_z
\quad\text{for every zero } \rho' \ne \rho.
\]
\item Every zero $\rho'$ of $\zeta$ satisfies
\[
|\Sig - \rho'| \;\ge\; D_{\min} \;>\; r',
\]
so $F$ is holomorphic on an open set containing the closed disk
$D(\Sig, r')$; moreover $252 - r' \ge 0.103 > 0$ and every point $w$ with
$|w - \Sig| = r'$ satisfies
\[
|w - \rho'| \;\ge\; \delta'
\quad\text{for every zero } \rho'.
\]
\end{enumerate}
\end{lemma}

\begin{proof}
(1) Let $\rho' = \beta' + i\gamma'$, $\rho' \ne \rho$. If
$\beta' \le \tfrac12$ then $|z - \rho'| \ge \Sig - \beta' \ge
\Sig - \tfrac12$. If $\beta' > \tfrac12$, then $\rho'$ lies in the slab
$\Sigma$ or not. If it does, (S1) forces $\rho' = \rho$ (there is only
one zero with $\re s > 1/2$ in the slab), contradiction. If not, then
--- since $\beta' > \tfrac12$ and $\rho' \notin \Sigma$, the defining
inequalities of $\Sigma$ force $\gamma' \le \gamma_1 - h$ or
$\gamma' \ge \gamma_2 + h$, i.e.\
$|\gamma' - \mathrm{mid}| \ge h + \Gamma/2 \ge h$ --- and then
\[
|z - \rho'|^2 = (\Sig - \beta')^2 + (\mathrm{mid} - \gamma')^2
\;\ge\; (\Sig - 1)^2 + (h + \Gamma/2)^2
\;\ge\; 251^2 + 17^2 > 251.5^2,
\]
using $\beta' < 1$ and (V4). Since $r_z < 251.5$, the disk claims
follow; $|z-\rho| \le \dbar < r_z$ is the definition of $\dbar$; and
$\re w \ge \Sig - r_z > \Sig - 251.5 = \tfrac12$ for $w$ in the
$z$-circle. The last displayed inequality of (1) then follows for
$\beta' \le \tfrac12$ from
$|w - \rho'| \ge \re w - \beta' \ge \Sig - r_z - \tfrac12 =
\varepsilon_z$, and for $\beta' > \tfrac12$ outside the slab from
\[
|w-\rho'| \ge |z - \rho'| - r_z > 251.5 - r_z = \varepsilon_z.
\]

(2) Let $\rho' = \beta' + i\gamma'$ be any zero. If $\beta' \le
\tfrac12$, then by Input~\ref{lem:zlow}(1),(2), $|\gamma'| \ge
14.1347$, so
\[
|\Sig - \rho'|^2 = (\Sig - \beta')^2 + \gamma'^2 \ge
(\Sig - \tfrac12)^2 + 14.1347^2 = D_{\mathrm{line}}^2.
\]
If $\beta' > \tfrac12$ and $|\gamma'| \le 30$, then
Input~\ref{lem:zlow}(3) gives $\beta' = \tfrac12$, impossible; so
$\beta' > \tfrac12$ forces $|\gamma'| > 30$ unless $\rho' = \rho$.
For such $\rho'$ (i.e.\ $\beta' > \tfrac12$, $\rho' \neq \rho$, hence
$|\gamma'| > 30$):
\[
|\Sig - \rho'|^2 \ge (\Sig - 1)^2 + 30^2 = 251^2 + 900
> D_{\mathrm{line}}^2 > D_{\min}^2.
\]
Finally if $\rho' = \rho$ then $\rho \in R$ gives
$|\Sig - \rho| \ge D_R \ge D_{\min}$. Hence every zero is at distance
$\ge D_{\min}$ from $\Sig$. The remaining poles of $F$ are at $s = 0$
(distance $252 > D_{\min} > r'$) and at the trivial zeros $-2, -4,
\ldots$ (distance $\ge 254$), so $F$ is holomorphic on
$\{|w - \Sig| < D_{\min}\}$, an open set containing the closed disk
$D(\Sig, r')$. Also $r' < D_{\min} \le D_{\mathrm{line}}$ gives
$252 - r' > 252 - 251.90 = 0.10$. For $w$ in the $\Sig$-circle and any
zero $\rho'$: $|w - \rho'| \ge |\Sig - \rho'| - r' \ge D_{\min} - r'
= \delta'$.
\end{proof}

\subsection{Explicit bounds on the two circles}

\begin{assumption}[local zero count; Backlund, conservative]\label{lem:backlund}
For every real $t$, the number $N_{\mathrm{loc}}(t)$ of zeros
$\rho = \beta + i\gamma$ of $\zeta$ with $|\gamma - t| \le 1$ satisfies
\[
N_{\mathrm{loc}}(t) \;\le\; N_b(t) := 0.6\log(|t| + 2) + 2.
\]
For $|t| \le 21.5$ one may take $N_{\mathrm{loc}}(t) \le 3$. (From
Backlund's explicit estimates for $N(T)$, as sharpened by
\cite{Tru14, Bac14}; the constants are deliberately generous.)
\end{assumption}

\begin{assumption}[partial-fraction bound for $\zeta'/\zeta$; standard,
conservative]\label{lem:zeta-prime-prime}
Let $s \in \C$ with $-1 \le \re s \le 3$, and let $\delta \ge 10^{-6}$
be such that every zero of $\zeta$ and the point $s=1$ are at distance
$\ge \delta$ from $s$. Then
\[
\Bigl|\frac{\zeta'}{\zeta}(s)\Bigr|
\;\le\; \frac{N_{\mathrm{loc}}(\im s) + 2}{\delta}
\;+\; 12\log(|\im s| + 4) + 12,
\]
and for $\re s \ge 3$,
$\bigl|\zeta'/\zeta(s)\bigr| \le
-\zeta'(3)/\zeta(3) \le 1$.
\end{assumption}

\noindent Lemmas~\ref{lem:backlund} and
\ref{lem:zeta-prime-prime} are quoted with deliberately generous
constants. The partial-fraction bound is the classical local
representation
\[
\frac{\zeta'}{\zeta}(s) \;=\;
\sum_{|\gamma - \im s| \le 1}\frac{1}{s-\rho}
\;+\; O\bigl(\log|\im s|\bigr)
\]
(Hadamard product for $\xi$, tail by partial
summation, $\Gamma$-factor by elementary digamma bounds; cf.\
Titchmarsh \cite{Tit51}, \S\S2.12, 3.5, and Montgomery--Vaughan
\cite{MV07}, \S10); the conservative constants here only ever enter our
bounds through explicit rational formulas, so any proven explicit
version is a drop-in replacement. The small-$|\im s|$ range is covered
by the same formula, since $N_{\mathrm{loc}}(t) = 0$ for $|t| \le 13$
(Input~\ref{lem:zlow}). A self-contained derivation is deferred to the
effective-shadow phase.

\begin{definition}[circle sup-bounds]\label{def:Mz}\
By (V5) all of $\varepsilon_z$, $\delta'$, $W_1$, $W_1'$ are $\ge
10^{-6}$; also $W_0 > \tfrac12$, $W_0' \ge 0.10$ and $W_1' > 0$ by
Lemma~\ref{lem:circles}(2) and $\dbar > 251$. Define
\begin{align*}
B_z &:= \frac{N_b(t_{\max}) + 2}{\min(\varepsilon_z,\, W_1)}
\;+\; 12\log(t_{\max} + 4) + 12,\\
B' &:= \frac{3 + 2}{\min(\delta',\, W_1')} + 12\log(26) + 12,
\end{align*}
and
\[
M_z := \frac{B_z}{W_0\,W_1} + \frac{1}{W_1^2} + \frac{1.6}{W_1}
\;+\; \frac{1}{M_1 M_2\, \varepsilon_z},
\qquad
M' := \frac{B'}{W_0'\,W_1'} + \frac{1}{(W_1')^2} + \frac{1.6}{W_1'} .
\]
\end{definition}

\begin{remark}[$M_z$ bounds $|h|$, not $|F|$]\label{rem:cancellation}\
\begin{enumerate}
\item $M_z$ is an explicit upper bound for
$|h| = |F - c_\rho/(s-\rho)|$ on the $z$-circle
(Lemma~\ref{lem:Mz-hold}(1)); the term
$(M_1M_2\varepsilon_z)^{-1}$ absorbs the principal part's boundary
contribution. $M'$ bounds $|F|$ on the $\Sig$-circle.
\item $F$ is holomorphic at $s=1$; our term-by-term bound treats the
pole of $\zeta'/\zeta$ at $1$ via the $W_1$-denominators, wasting the
cancellation when a circle passes near $s=1$. This only affects the
size of constants (a larger $M_z$ or $M'$ merely raises the explicit
thresholds $j_1$ and $J_0$ below), never correctness or computability.
\item $W_1 = ||z-1|-r_z|$ could in principle vanish for an arithmetic
accident of the certificate data; (V5) excludes this by a trivial
rational check (compare rational squares), and the same check
guarantees $W_1 \ge 10^{-6}$.
\end{enumerate}
\end{remark}

\begin{lemma}[sup-bounds hold]\label{lem:Mz-hold}\
\begin{enumerate}
\item For every $w$ with $|w - z| = r_z$:
$|w| \ge W_0 = \tfrac12$, $|w - 1| \ge W_1 \ge 10^{-6}$, and
\[
\Bigl|F(w) - \frac{c_\rho}{w-\rho}\Bigr| \;\le\; M_z.
\]
\item For every $w$ with $|w - \Sig| = r'$:
$|w| \ge W_0' \ge 0.10$, $|w-1| \ge W_1' \ge 10^{-6}$, and
\[
|F(w)| \;\le\; M'.
\]
\end{enumerate}
\end{lemma}

\begin{proof}
(1) Since $|z| \ge \Sig = 252$ and $r_z < 251.5$,
$|w| \ge |z| - r_z \ge 252 - 251.5 = \tfrac12 = W_0$.
For any fixed point $p$, the nearest point of the circle
$\{|w-z| = r_z\}$ to $p$ is at distance $\bigl||z-p| - r_z\bigr|$
(along the segment from $z$ through $p$), so
$|w-1| \ge \bigl||z-1|-r_z\bigr| = W_1 \ge 10^{-6}$, the last step by
(V5). By Lemma~\ref{lem:circles}(1), every zero $\rho' \ne \rho$ is at
distance $\ge \varepsilon_z \ge 10^{-6}$ from $w$, and $\rho$ itself at
distance $\ge r_z - \dbar = \varepsilon_z$ from $w$. Thus every zero
and the point $1$ are at distance $\ge \delta := \min(\varepsilon_z,
W_1) \ge 10^{-6}$ from $w$. If $\re w \le 3$,
Input~\ref{lem:zeta-prime-prime} gives
\begin{align*}
|\zeta'/\zeta(w)| \;&\le\; \frac{N_b(|\im w|)+2}{\delta}
+ 12\log(|\im w|+4) + 12\\
\;&\le\; \frac{N_b(t_{\max})+2}{\min(\varepsilon_z, W_1)}
+ 12\log(t_{\max}+4) + 12 \;=\; B_z,
\end{align*}
where the second inequality uses $|\im w| \le t_{\max}$
(\S\ref{sec:derived}) and the monotonicity of $N_b$ and $\log$. If
$\re w > 3$, then $|\zeta'/\zeta(w)| \le 1 \le B_z$ by
Input~\ref{lem:zeta-prime-prime} (second clause), since $B_z \ge 12 > 1$.
From \eqref{eq:Fdef},
\[
|F(w)| \le \frac{B_z}{W_0 W_1} + \frac{1}{W_1^2} + \frac{1.6}{W_1},
\]
and, since $|w - \rho| \ge \varepsilon_z$ and
$|c_\rho| \le 1/(M_1M_2)$,
\[
\Bigl|F(w) - \frac{c_\rho}{w-\rho}\Bigr|
\;\le\; |F(w)| + \frac{1}{M_1M_2\,\varepsilon_z} \;\le\; M_z .
\]

(2) $|w| \ge \Sig - r' = W_0' \ge 0.10$; $|w-1| \ge
\bigl||\Sig-1|-r'\bigr| = |251 - r'| = r' - 251 = W_1'$ (as $r' > 251$),
which is $\ge 10^{-6}$ by (V5). Every zero is at distance $\ge \delta' \ge 10^{-6}$ from $w$
(Lemma~\ref{lem:circles}(2)), and $|\im w| \le \sqrt{(r')^2 - 251^2} <
\sqrt{251.9^2 - 251^2} < 21.5$, so $N_{\mathrm{loc}}(\im w) \le 3$
(Input~\ref{lem:backlund}, second clause). If $\re w \le 3$,
Input~\ref{lem:zeta-prime-prime} with $\delta = \min(\delta', W_1')$
gives $|\zeta'/\zeta(w)| \le B'$; if $\re w > 3$,
$|\zeta'/\zeta(w)| \le 1 \le B'$. Hence $|F(w)| \le B'/(W_0'W_1') +
1/(W_1')^2 + 1.6/W_1' = M'$.
\end{proof}

\begin{remark}
The proof of Lemma~\ref{lem:pole-lower} is the effectivized home of
Rubinstein's existential step (``choose $t$ so that $z$ has a unique
nearest pole, at distance $d < R$; choose $d < r < R$ so the disk of
radius $r$ about $z$ contains just that pole''). In our version the
choice $z = \Sig + i\,\mathrm{mid}$ is explicit, and the role of the
``unique nearest pole'' is played by the certificate (S1): the disk
$D(z, r_z)$ contains the certified zero and, by the slab, no other
zero to the right of the line, while everything on or left of the line
is geometrically farther than $r_z$. No other zero can interfere with
the principal part.
\end{remark}

\subsection{The pole lower bound}

\begin{lemma}[pole lower bound]\label{lem:pole-lower}
Let the certificate be valid (Definition~\ref{def:cert}) and let (S1)
hold. Put
\[
j_1 := \left\lceil
\frac{\log\bigl(2 M_1 M_2\, M_z\, \dbar\bigr)}
{\log(r_z/\dbar)} \right\rceil .
\]
Then for every integer $j \ge j_1$,
\[
|A_j(z)| \;\ge\; \frac{1}{2 M_1 M_2}\;\dbar^{-j-1}.
\]
\end{lemma}

\begin{proof}
Let $q(s) := F(s) - c_\rho/(s-\rho)$ (we avoid the letter $h$, reserved
for the slab half-width). By Lemma~\ref{lem:circles}(1), the only zero
of $\zeta$ in the open disk $\{w : |w - z| < \Sig - \tfrac12\}$ is
$\rho$; hence $q$ is holomorphic on that disk, which contains the
closed disk $D(z, r_z)$ because $r_z < \Sig - \tfrac12$.

For $|u| < d := |z - \rho|$ (note $d \le \dbar$), the decomposition
\eqref{eq:Fdef} minus the principal part gives
\[
F(z+u) \;=\; \frac{c_\rho}{z+u-\rho} + q(z+u)
\;=\; \sum_{k \ge 0} c_\rho (z-\rho)^{-k-1} (-u)^k \;+\; q(z+u),
\]
using $\frac{1}{z+u-\rho} = \frac{1}{z-\rho}\cdot
\frac{1}{1 + u/(z-\rho)}$. Reading off Taylor coefficients at $u = 0$
and applying the sign $(-1)^j$ in the definition of $A_j$,
\[
A_j(z) \;=\; c_\rho (z-\rho)^{-j-1} \;+\; (-1)^j [u^j]\,q(z+u).
\]
By Cauchy's coefficient estimate on the circle $|u| = r_z$, where
$|q| \le M_z$ by Lemma~\ref{lem:Mz-hold}(1),
$|[u^j]q(z+u)| \le M_z\, r_z^{-j}$. Using
\eqref{eq:residue} and $d \le \dbar$,
\[
|A_j(z)| \;\ge\; \frac{\dbar^{-j-1}}{M_1 M_2} \;-\; M_z\, r_z^{-j}
\;=\; \dbar^{-j-1}\Bigl(\frac{1}{M_1M_2} \;-\; \dbar\, M_z
\bigl(\tfrac{\dbar}{r_z}\bigr)^{j}\Bigr),
\]
and the parenthesis is $\ge \frac{1}{2M_1M_2}$ once
$2 M_1 M_2\, \dbar\, M_z \bigl(\frac{\dbar}{r_z}\bigr)^{j} \le 1$,
i.e.\ once $j \ge j_1$ (note $\log(r_z/\dbar) > 0$ since $r_z > \dbar$).
\end{proof}

\section{The main theorem}\label{sec:main}

\begin{definition}[thresholds and the margin]\label{def:thresholds}
Given a valid certificate (Definition~\ref{def:cert}), define the
following further computable quantities:
\begin{align*}
\mu &:= \frac{\Sig-\theta}{r'}, \qquad
j_2 := \max\Bigl(0,\;
\Bigl\lceil \frac{\log\bigl(8 M_1 M_2\, M'\,(\Sig-\theta)/\lambda\bigr)}
{\log(r'/\dbar)} \Bigr\rceil\Bigr),\\
J_0 &:= \left\lceil
\frac{\log\bigl(\tfrac{16}{3} M_1 M_2\, (C + 54)\bigr)}{\log\lambda}
\right\rceil,
\qquad
J_* := \max(j_1, j_2, J_0),\\
U_0 &:= \max\Bigl(3,\; \frac{J_* + \tfrac12}{307.5}\Bigr),
\qquad
Y_0 := e^{U_0},
\end{align*}
and, for integers $j \ge \max(j_1, j_2)$, the \emph{margin function}
\[
Q(j) \;:=\; \frac{3\,\lambda^{j+1}}{16\,M_1 M_2} \;-\; 52.5 .
\]
\end{definition}

\begin{theorem}[T1: effective windowed Landau certificate theorem]
\label{thm:main}
Let
$(\theta, C, \beta_1, \beta_2, \gamma_1, \gamma_2, h)$
be a valid certificate, i.e.\ satisfying the validity conditions
{\rm (V1)--(V5)} of Definition~\ref{def:cert} and the certified
statement {\rm (S1)}. Let $Y \ge Y_0$ be rational and put
$U := \log Y$, and let $j$ be the nearest integer to $1.23(\Sig-2)U
= 307.5\,U$. Then:
\begin{enumerate}
\item $j \ge J_*$, and consequently $Q(j) \ge C + 1$; put
$m := Q(j) - C \ge 1$.
\item The ceiling \emph{fails} on the window: there exists
$x \in [Y, Y^{3/2}]$ with
\[
G(x) \;\ge\; Q(j)\, x^{\theta} \;\ge\; (C + m)\,x^{\theta}.
\]
\item Moreover such an $x$ can be found by a terminating algorithm:
the grid search of \S\ref{sec:algorithm}
(Proposition~\ref{prop:search}) outputs a rational
$x \in [Y, Y^{3/2}]$ with
\[
G(x) \;\ge\; \bigl(C + \tfrac{m}{2}\bigr)\,x^{\theta},
\]
together with a finite verification of this inequality (a sum over
the prime powers $\le Y^{3/2}$ evaluated with rigorous rational
interval arithmetic). The computation terminates in a number of
steps that is bounded by an explicit polynomial in $Y$ and in the
bit-lengths of the input data.
\end{enumerate}
\end{theorem}

\begin{proof}
(1) Since $U \ge U_0 \ge 3$, Lemma~\ref{lem:j-range} applies; since
also $U \ge (J_* + \tfrac12)/307.5$, we get
$j \ge 1.23 \cdot 250\, U - \tfrac12 \ge J_* \ge \max(j_1, j_2, J_0)$.
By construction $Q$ is increasing in $j$ (as $\lambda > 1$), and at
$j = J_0$,
\[
\begin{split}
\lambda^{J_0} \;&\ge\; \tfrac{16}{3} M_1 M_2 (C+54)\\
\;&\Longrightarrow\;
Q(J_0) \;\ge\; \tfrac{3}{16M_1M_2}\cdot\tfrac{16}{3}M_1M_2(C+54)
- 52.5 \;=\; C + 1.5,
\end{split}
\]
so $Q(j) \ge Q(J_0) \ge C + 1.5$ and $m = Q(j) - C \ge 1.5$.

(2) Suppose, for contradiction, that $G(x) < Q(j) x^{\theta}$ for all
$x \in [Y, Y^{3/2}]$. The function $x \mapsto G(x)/x^{\theta}$ is
continuous on the compact window $[Y, Y^{3/2}]$ (Lemma~\ref{lem:G-size}
(2); $x^\theta$ is continuous and positive), so it attains its maximum
$C'$ there, and $C' < Q(j)$. Then
\[
G(x) \;\le\; C' x^{\theta} \quad\text{for all } x \in [Y, Y^{3/2}]
\]
with $C' \ge 0$. Apply Lemma~\ref{lem:onesided} (with ceiling constant
$C'$) and Corollary~\ref{cor:upper}: the latter is applicable because
$F$ is holomorphic on an open set containing the closed disk
$D(\Sig, r')$ by Lemma~\ref{lem:circles}(2), so
Lemma~\ref{lem:sigcauchy} gives
\[
|A_j(z)| \;\le\; (2C' + 105)(\Sig-\theta)^{-j-1} + M' (r')^{-j}
\qquad (z = \Sig + i\,\mathrm{mid}).
\]
On the other hand, $j \ge j_1$, so Lemma~\ref{lem:pole-lower} gives
$|A_j(z)| \ge \dbar^{-j-1}/(2M_1M_2)$. Combining and multiplying by
$(\Sig-\theta)^{j+1}$:
\[
\frac{\lambda^{j+1}}{2M_1M_2}
\;\le\; 2C' + 105 + M'(\Sig-\theta)\,\mu^{j}.
\]
Since $j \ge j_2$ and $\lambda/\mu = r'/\dbar$, the last term obeys
$M'(\Sig-\theta)\mu^j \le \lambda^{j+1}/(8M_1M_2)$ by the definition
of $j_2$. Hence
\[
\frac{3\lambda^{j+1}}{8M_1M_2} \;\le\; 2C' + 105,
\qquad\text{i.e.}\qquad
C' \;\ge\; \frac{3\lambda^{j+1}}{16M_1M_2} - 52.5 \;=\; Q(j),
\]
contradicting $C' < Q(j)$. So the maximum of $G(x)/x^{\theta}$ over
the window is $\ge Q(j)$, attained at some $x^* \in [Y, Y^{3/2}]$,
which proves (2).

(3) is Proposition~\ref{prop:search} below.
\end{proof}

\begin{corollary}[windowed Landau, certificate form]\label{cor:landau}
Fix $\theta \in [\tfrac12, 1) \cap \Q$ and $C \ge 1$. If for some
rational $Y \ge Y_0$ the ceiling
\[
G(x) \;\le\; C\,x^{\theta} \qquad\text{for all } x \in [Y, Y^{3/2}]
\]
holds (whether or not anyone verifies it), then no valid certificate
with this $\theta$ exists: there is no rectangle $R$ with
$\beta_1 > \theta$ and no slab for which the certified statement (S1)
is true. In particular, under the additional assumption that every
zero of $\zeta$ with real part $> \theta$ is simple and admits a
valid certificate (i.e.\ is isolated from all other off-critical-line
zeros by a slab of half-height $\ge \tfrac{\Gamma}{2} + 17$; see
\S\ref{sec:limits}), $\zeta$ has no zero with real part $> \theta$.
\end{corollary}

\begin{proof}
Immediate from Theorem~\ref{thm:main}(2): a valid certificate would
force the inequality $G(x^*) > Cx^{*\theta}$ at some point
$x^* \in [Y, Y^{3/2}]$ (indeed $G(x^*) \ge (C+1)x^{*\theta}$),
contradicting the ceiling.
\end{proof}

\begin{remark}[relation to Lemma~\ref{lem:rubinstein}]
If the ceiling $G \le x^{\theta}$ (say $C = 1$) holds on windows
$[Y_k, Y_k^{3/2}]$ with $Y_k \to \infty$, then for every fixed
certificate the ceiling eventually contradicts it, i.e.\ no valid
certificate exists, and if off-line zeros are certifiable as above
this is Rubinstein's windowed Landau lemma. The gap between this
corollary and the unconditional classical statement is exactly the
cluster issue discussed in \S\ref{sec:limits}: the present argument
requires the certified zero to be the unique off-line zero in its
slab, which is exactly what the classical existential argument does
\emph{not} require. We insist on the distinction: what is made
effective here is the certificate-level statement, which is the form
the ACL2 checker (target T3) consumes; the cluster-free general
statement is a consequence under the stated (standard, checkable)
non-clustering hypothesis.
\end{remark}

\begin{proposition}[the address search]\label{prop:search}
In the setting of Theorem~\ref{thm:main}, define
\[
L_h := L_G(Y^{3/2}) + C + \tfrac{m}{2},
\qquad
\delta_{\mathrm{grid}} := \frac{m\,Y^{\theta}}{4\,L_h},
\qquad
N_{\mathrm{grid}} := \left\lceil
\frac{Y^{3/2} - Y}{\delta_{\mathrm{grid}}} \right\rceil,
\]
where $L_h$ is a Lipschitz constant, on the window, for
$x \mapsto G(x) - \bigl(C + \tfrac{m}{2}\bigr)x^{\theta}$; and put
\[
x_k := \min\bigl(Y + k\,\delta_{\mathrm{grid}},\; Y^{3/2}\bigr)
\qquad\text{for } k = 0, 1, \ldots, N_{\mathrm{grid}} .
\]
Let $x^*$ be as in Theorem~\ref{thm:main}(2). Then some $k \le
N_{\mathrm{grid}}$ satisfies
\[
G(x_k) \;-\; \bigl(C + \tfrac{m}{2}\bigr)\,x_k^{\theta}
\;\ge\; \frac{m}{4}\,Y^{\theta} \;>\; 0,
\]
and for that $k$ (and in fact for every $k$) the left-hand side can be
rigorously certified, from above and below, by evaluating $G(x_k)$
as a finite sum over prime powers $n \le x_k$ with each of
$\log p$ ($p \le x_k$), $\log x_k$, $\Euler$, and $x_k^{\theta}$
computed by rational interval arithmetic of explicit width
$w$ given in the proof. Consequently the search over
$k = 0, \ldots, N_{\mathrm{grid}}$ terminates and outputs an address
$x_k$ with a certified
\[
G(x_k) \;\ge\; \bigl(C + \tfrac{m}{2}\bigr)\,x_k^{\theta}.
\]
\end{proposition}

\begin{proof}
Write $H(x) := G(x) - (C + \tfrac{m}{2})x^{\theta}$. On $[Y, Y^{3/2}]$,
$G$ is $L_G(Y^{3/2})$-Lipschitz (Lemma~\ref{lem:G-size}(2)) and
$x \mapsto x^{\theta}$ has $|\theta x^{\theta-1}| \le 1$ for
$x \ge Y \ge 1$ (as $0 < \theta < 1$), so $H$ is $L_h$-Lipschitz. At
$x^*$,
\[
H(x^*) \;=\; G(x^*) - \bigl(C+\tfrac{m}{2}\bigr) x^{*\theta}
\;\ge\; (C+m)x^{*\theta} - \bigl(C+\tfrac{m}{2}\bigr)x^{*\theta}
\;=\; \tfrac{m}{2}\,x^{*\theta} \;\ge\; \tfrac{m}{2}Y^{\theta}.
\]
The grid points satisfy $\max_k |x_k - x^*| \le \delta_{\mathrm{grid}}$
(the grid starts at $Y$, ends at $Y^{3/2}$, with spacing at most
$\delta_{\mathrm{grid}}$), so for the nearest grid point $x_k$ to
$x^*$:
\[
H(x_k) \;\ge\; H(x^*) - L_h\,\delta_{\mathrm{grid}}
\;\ge\; \tfrac{m}{2}Y^{\theta} - \tfrac{m}{4}Y^{\theta}
\;=\; \tfrac{m}{4}Y^{\theta} \;>\; 0.
\]
It remains to see that $H(x_k)$ is finitely certifiable. Fix $k$ and
write the finite sum
\[
G(x_k) \;=\; \sum_{\substack{n = p^a \le x_k}} (\log p)
\Bigl(\frac{x_k}{n} - 1\Bigr) \;-\; x_k \log x_k \;+\; (1+\Euler)\,x_k ,
\]
over the finitely many prime powers $n \le x_k$ (enumerated by a
sieve up to $\lfloor x_k \rfloor$, or up to $\lfloor Y^{3/2}\rfloor$
once and for all). Suppose each of the quantities
$\log p$ ($p \le x_k$), $\log x_k$ and $\Euler$ is computed by
rigorous interval arithmetic to width $w$, and $x_k^{\theta}$ to
width $2w$ (from $x_k^{\theta} = e^{\theta\log x_k}$). Using
$S(x_k) = \sum_{n \le x_k}\Lam(n)(x_k/n - 1) \le
\tfrac12 x_k(\log x_k)^2$
(Lemma~\ref{lem:G-size}(1) and its proof) and $n \le x_k$
throughout, the total propagated error in $G(x_k)$ is at most
\[
w\Bigl(\tfrac12 x_k(\log x_k)^2 + 2x_k + 1\Bigr) \;\le\;
w\,E_0, \qquad
E_0 := Y^{3/2}\bigl(\tfrac{1}{2}\bigl(\tfrac32\log Y\bigr)^2 + 3\bigr),
\]
and the error in $(C + \tfrac{m}{2})x_k^{\theta}$ is at most
$2w(C + \tfrac{m}{2})\,x_k^{\theta} \le 2w(C + \tfrac{m}{2})\,Y^{
3\theta/2}$. Choose
\[
w := \frac{m\,Y^{\theta}}{8\bigl(E_0 + 2(C+\tfrac{m}{2})Y^{3\theta/2}
+ 1\bigr)} \;>\; 0,
\]
a rational. Then the certified value of $H(x_k)$ differs from the true
one by at most $\tfrac{m}{8}Y^{\theta}$, so the good $k$ found above
certifies $H(x_k) \ge \tfrac{m}{4}Y^{\theta} - \tfrac{m}{8}Y^{\theta}
= \tfrac{m}{8}Y^{\theta} > 0$. Every quantity involved
($\log$ of a rational, $e^x$ of a rational, the sieve) is computable
to width $w$ by terminating rational procedures (e.g.\ alternating
series with explicit remainder bounds, or interval Newton
iterations), so each evaluation terminates, and there are
$N_{\mathrm{grid}} + 1$ of them. Finally,
\[
\begin{split}
N_{\mathrm{grid}} \;&\le\; \frac{4\,Y^{3/2}\,L_h}{m\,Y^{\theta}}
\;\le\; 4\,Y^{3/2 - \theta}\Bigl(\tfrac{L_G(Y^{3/2}) + C}{m}
+ \tfrac12\Bigr)\\
\;&\le\; 4\,Y^{3/2-\theta}\bigl(L_G(Y^{3/2}) + C + \tfrac12\bigr),
\end{split}
\]
using $m \ge 1$; together with the sieve cost (polynomial in
$Y^{3/2}$) and the per-point interval arithmetic (polynomial in
$\log(1/w)$ and $Y^{3/2}$), the total is bounded by an explicit
polynomial in $Y$ and the input bit-lengths.
\end{proof}

\begin{remark}[what the verifier actually checks]
The output of the search is \emph{self-certifying}: the pair
$(x_k, m/2)$ comes with an interval-arithmetic transcript from which
the inequality $G(x_k) \ge (C + \tfrac{m}{2})x_k^{\theta}$ follows by
finitely many rational comparisons. Theorem~\ref{thm:main} guarantees
that the search never runs to exhaustion on valid input: reaching
$k = N_{\mathrm{grid}}$ without a certified violation certifies that
the input was \emph{not} valid --- either some arithmetic check
failed silently, or (S1) is false, i.e.\ the upstream zero-isolation
certificate is bogus. This is precisely the shape required by the
ACL2 checker (target T3): the analytic content is concentrated in the
hypotheses, and the verification is a finite computation.
\end{remark}

\section{The certificate algorithm}\label{sec:algorithm}

We specify the algorithm of Proposition~\ref{prop:search} as an
implementable procedure. A machine-readable specification with exact
input/output formats is the companion file
\texttt{classical/certificates.md}; the two are kept in sync.

\begin{quote}
\textbf{INPUT:} rationals
$\theta, C, Y, \beta_1, \beta_2, \gamma_1, \gamma_2, h$
(together with the upstream certification of (S1), which the
algorithm treats as a hypothesis and never checks).\\[4pt]
\textbf{OUTPUT:} either \textsc{invalid} (with the first failed
check named), or a pair $(x, m/2)$ with $x \in \Q \cap [Y, Y^{3/2}]$
and a verification transcript establishing
$G(x) \ge (C + \tfrac{m}{2})x^{\theta}$.
\end{quote}

\begin{enumerate}
\item \emph{Data ranges.} Check the ranges $\theta \in [\tfrac12, 1)$,
$C \ge 1$, $Y \ge 20$; the orderings $\beta_1 < \beta_2$,
$\gamma_1 < \gamma_2$; the sign $h > 0$; compute $\Gamma$,
$\mathrm{mid}$.
\item \emph{Derived quantities.} Compute certified rational intervals
for the fourteen quantities of \S\ref{sec:derived}, namely
\[
\dbar,\quad M_1,\quad M_2,\quad D_R,\quad D_{\min},\quad \lambda,\quad
r_z,\quad \varepsilon_z,\quad W_1,
\]
\[
r',\quad \delta',\quad W_0',\quad W_1',\quad t_{\max},
\]
each to enough precision that the subsequent comparisons are decided.
\item \emph{Validity.} Check (V1), (V2), (V3), (V4), (V5).
\item \emph{Circle bounds.} Compute $B_z, B', M_z, M'$
(Definition~\ref{def:Mz}).
\item \emph{Thresholds.} Compute $j_1, j_2, J_0, J_*, U_0, Y_0$
(Definitions~\ref{def:thresholds} and \ref{lem:pole-lower});
check $Y \ge Y_0$. Compute $U := \log Y$ to within $10^{-3}$ and
$j :=$ the nearest integer to $307.5\,U$.
\item \emph{Margin.} Compute a certified lower bound for
$Q(j)$; check $Q(j) > C + 1$; put $m := Q(j) - C$.
\item \emph{Grid and precision.} Compute $L_h$,
$\delta_{\mathrm{grid}}$, $N_{\mathrm{grid}}$, $w$
(Proposition~\ref{prop:search}).
\item \emph{Sieve.} Enumerate the prime powers $\le Y^{3/2}$.
\item \emph{Search.} For $k = 0, \ldots, N_{\mathrm{grid}}$:
evaluate $G(x_k)$ by interval arithmetic of width $w$; if the
certified value of $H(x_k)$ is positive, output
$(x_k,\, m/2)$ and the transcript, and halt.
\item \emph{Failure.} If the loop completes without success, output
\textsc{invalid}: by Theorem~\ref{thm:main} some input condition
failed (most plausibly (S1), i.e.\ the upstream certificate).
\end{enumerate}

\begin{remark}[complexity]
The dominant costs are the sieve to $Y^{3/2}$ and the
$N_{\mathrm{grid}} \le 4\,Y^{3/2-\theta}(L_G(Y^{3/2}) + C +
\tfrac12)$ interval evaluations, each of polynomial cost; the whole
procedure is bounded by an explicit polynomial in $Y$ and the input
bit-lengths (Proposition~\ref{prop:search}). No step involves limits,
infinite sums, or non-constructive choices: every real quantity is
handled through rational intervals with explicitly computed widths.
\end{remark}

\section{Limitations, variants, and notes for the next phases}\label{sec:limits}

\subsection{What the certificate excludes: clusters}\label{sec:clusters}

The certified statement (S1) demands that the slab
$\Sigma = \{\re s > \tfrac12\} \times (\gamma_1 - h, \gamma_2 + h)$
--- which by (V4) has half-height at least $\Gamma/2 + 17 \approx 17$
--- contain \emph{exactly one} zero off the critical line. Two off-line
zeros whose heights differ by less than about $34$ therefore cannot be
handled by the present argument, either individually or jointly, and
the certificate cannot be issued for either. Consequently:

\begin{itemize}
\item Corollary~\ref{cor:landau}, in its unconditional form, needs the
hypothesis that every off-line zero is simple and admits a valid
certificate --- a genuine restriction relative to the classical
statement (Lemma~\ref{lem:rubinstein}), whose existential choice of the
evaluation point (``choose $t$ so that $z = \Sig + it$ has a unique
nearest pole'') absorbs clusters automatically.
\item The classical nearest-pole choice has not been effectivized by
the fixed-slab argument above. This is a limitation of the present
certificate, not an impossibility of finite separation data: a complete
certified finite cover can carry explicit separation information.
The repository now verifies rational tie-avoiding heights, nearest-point
selection, and isolating circles for finite rational point data in ACL2.
A box-circle checker has a Lean interpretation for actual complex zeros,
conditional on a complete cover and uniqueness in the selected box.
Generating and refining suitable covers for arbitrary zeros, and
deriving the required analytic estimates from them, remain proof
obligations. Until those are discharged, the single-zero certificate
does not establish the full windowed Landau reduction in arithmetic.
\end{itemize}

Two further routes may extend the reach; neither is established here:

\begin{enumerate}
\item \emph{Phase-averaging.} With several poles $\rho_1, \ldots,
\rho_m$ in the disk, subtract all principal parts and note that a
conjugate pair contributes $2\re\bigl(c_\rho (z-\rho)^{-j-1}\bigr)$,
whose phase precesses per unit increment of $j$ by the amount
$2\arg(z-\rho) \ge 2\arctan\bigl(|\im\rho|/(\Sig-\beta)\bigr)$. The
available range of $j$ (the kernel window, about $9.4\,U$ consecutive
integers) contains a favorable phase whenever
$2\pi/\arg(z-\rho)$ is explicit and small relative to it. This works
for pairs bounded away from the real axis by an explicit margin ---
which could be added as a certificate condition --- but not
uniformly, since a pair can sit at the evaluation height.
\item \emph{Structured isolation certificates.} Let the certificate
describe the cluster: a rectangle containing $m \ge 1$ simple zeros
with certified pairwise separations, all other zeros excluded by the
slab. The lower bound becomes a comparison between an explicit
$m$-term exponential sum and the ceiling term; with certified
separations the dominant pole wins at an explicit threshold. This is
routine but constant-heavy, and is planned as the Phase 1.5 companion.
\end{enumerate}

A multiple zero at a single point is \emph{not} an obstruction of this
kind: the pole of $F$ is then still simple, with residue
$-\nu_\rho/(\rho(\rho-1))$, and $|c_\rho| \ge \nu_\rho/(M_1M_2) \ge
1/(M_1M_2)$, so the argument goes through verbatim once (S1) is
phrased with multiplicity (as it is: ``exactly one zero, counted with
multiplicity'' demands simplicity only when the count is $1$; the
reader should mentally allow ``exactly one zero of multiplicity
$\nu$'' with the residue bound scaled by $\nu$).

\subsection{The Dirichlet $L$-function variant}\label{sec:lfunction}

The same argument runs for Dirichlet $L$-functions --- the form the
surrounding architecture consumes on its automorphic side. The
archive's \texttt{Gchi}/\texttt{Schi} live in
\texttt{StatementOC.lean}; \texttt{docs/statement.md} \S4 records the
matching statements. We record the shape here; full
constant-tracking is planned as a companion note (target T1-L).

For a Dirichlet character $\chi \pmod q$, define the real statistic
\[
S_\chi(x) := \sum_{n \le x} \Lam(n)\,\re\chi(n)\Bigl(\frac{x}{n} - 1\Bigr),
\]
and
\[
G_\chi(x) := S_\chi(x) + \re\frac{L'}{L}(1, \chi)\cdot x .
\]
(for $\chi = 1$ this is not needed; the principal character reduces to
$\zeta$). Its Mellin transform is, for $\re s > 1$,
\[
F_\chi(s) = \frac{1}{s(s-1)}\cdot\Bigl(-\frac{L'}{L}(s,\chi)
- \frac{\bar L'}{\bar L}(s,\bar\chi)\Bigr)\Big/ 2
+ (\text{elementary correction}),
\]
a meromorphic function whose poles in the critical strip are exactly
the zeros of $L(\cdot,\chi)$ and of $L(\cdot,\bar\chi)$ (conjugate
zeros), with residue $-\nu_\rho/(\rho(\rho-1))$ as before. The proof
of Theorem~\ref{thm:main} then goes through with:
\begin{itemize}
\item the classical inputs (Input~\ref{lem:zlow},
\ref{lem:backlund}, \ref{lem:zeta-prime-prime}) replaced by their
per-conductor analogues, supplied as certificates by upstream
computation (verified GRH-type data for small conductors exists in the
literature; for arbitrary conductors this is additional upstream work,
exactly as for $\zeta$);
\item real zeros (the Siegel-zero scenario) handled \emph{without
change} by the rectangle geometry --- the certificate places them in
$R$ with $\beta_1 > \theta$, and nothing in \S\ref{sec:certificate}
uses non-reality of the certified zero;
\item the same absolute constants ($\Sig = 252$, $c = 0.008$,
$\eta_*$, window exponent $3/2$), since none of them touched $\zeta$
except through the quoted inputs.
\end{itemize}

\subsection{Notes toward the Lean two-layer and the PA track}\label{sec:pa}

The Lean formalization (target T2) is planned in two layers
(\texttt{../docs/roadmap.md}, Phase 2): a \emph{soft layer} in
ordinary mathlib style, and a \emph{choice-free effective shadow}
interpretable in first-order arithmetic, feeding Phase 3 (reverse-math
placement) and Phase 4 (an ACL2 certificate checker). The present
paper was written to make both layers mechanical, and we record the
map:

\begin{itemize}
\item \emph{Hypotheses, not theorems.} Every genuinely analytic input
--- (S1), Input~\ref{lem:zlow} (low zero data),
Input~\ref{lem:backlund} (local zero count),
Input~\ref{lem:zeta-prime-prime} (the partial-fraction bound) --- enters
as a hypothesis of arithmetic form: bounded-quantifier statements about
values of $\zeta$ (equivalently, about terminating rational interval
computations with explicit outputs). In the shadow they are never
proved; they are consumed. This is the hypothesis-injection discipline,
and it is what keeps the PA-track low: the transcendental content is
quarantined in the certificate.
\item \emph{The soft layer} proves: Lemma~\ref{lem:rubinstein} (the
classical statement, matching \cite{Rub26}), and
Theorem~\ref{thm:main} over the reals with (S1) and the quoted inputs
as hypotheses. Ordinary classical analysis is allowed internally; the
target is correctness of the statement and proof structure.
\item \emph{The shadow} re-proves Theorem~\ref{thm:main} with: reals
appearing only through comparisons $x \le q$ with explicit rationals
$q$ and explicit error bounds; every existential replaced by the
computed witness (the grid point $x_k$; the thresholds $j_1, j_2,
J_0$; the precisions); the Mellin integral replaced by finite partial
sums with explicit remainders (the ``truncated kernel'' variant the
roadmap anticipated: the tail Lemmas~\ref{lem:tail}, \ref{lem:tail2}
are precisely the remainders needed); the Cauchy coefficient bounds
replaced by interval quadrature with explicit subdivisions (standard
effectivization of Cauchy's estimate for functions given with
certified boundary sup-bounds). No limit remains in the statement
path: the ``for all large $j$'' of the classical proof has been
replaced by the computed integer thresholds.
\item \emph{Expected strength.} The window conversion itself (kernel
inequalities, grid search, interval arithmetic) is elementary
recursive arithmetic; we expect the shadow to land in
$\mathrm{RCA}_0$-to-$\mathrm{I}\Sigma_1$ territory \emph{given} the
hypotheses, with the hypotheses themselves carrying whatever strength
the analytic facts have. Deciding that precisely is Phase 3's
deliverable; the honest expectation, and its reason, is recorded here
so that Phase 3 inherits a target, not a hope.
\item \emph{Statement fingerprint.} The shadow's main statement, in
arithmetic form, is:
``for all rational tuples satisfying (V1)--(V5), if (S1) and the
quoted inputs hold in their arithmetic forms, then the search of
\S\ref{sec:algorithm} returns $k \le N_{\mathrm{grid}}$ with
$\mathrm{check}(x_k)$ true'' --- a $\Pi_2$-shaped statement with all
witnesses computed; the ACL2 checker verifies exactly the
$\mathrm{check}$ part plus the threshold arithmetic.
\end{itemize}

\section{The constant registry}\label{sec:registry}

Every constant below is rational (or an explicitly computed rational
bound for a transcendental quantity, used only through certified
rational intervals). Irrational constants ($e$, $\log$, $\pi$,
$\Euler$) appear in the proof only inside such intervals.

{\footnotesize
\begin{longtable}{@{}p{0.23\linewidth}p{0.39\linewidth}p{0.30\linewidth}@{}}
\hline
\textbf{Symbol} & \textbf{Value / definition} & \textbf{Role}\\
\hline
\endhead
$\Sig$ & $252$ & abscissa of evaluation; least value with
$(\Sig-\theta)/(\Sig-2) \le e^{c}$ for all $\theta \in (0,1)$\\
$c$ & $0.008$ & tail decay; pays for the $(\Sig-2) \to
(\Sig-\theta)$ scale conversion\\
$\eta_*$ & $0.983$ & cap on the kernel outside the window
($2\eta_*/(1+\eta_*) < e^{-c}$)\\
$[v_-, v_+]$ & $[0.88, 1.12]$ & fixed interval about the kernel peak;
$b \ge 0.9921$ on it\\
$b(v) = ve^{1-v}$ & --- & the P\'olya kernel profile\\
$1.23$, $307.5$ & $1.23(\Sig-2) = 307.5$ & the $j$-choice:
$j = \lfloor 307.5\,U \rceil$; valid range $[1.2122, 1.25](\Sig-2)U$\\
$C_0$ & $52 = 4 \times 13$ & tail constant for
$\int_E 4e^{2u}K_j$, $E$ outside the window\\
$C_1$ & $105$ & the additive ceiling constant in
\eqref{eq:onesided}\\
$4$ & $|G(x)| \le 4x^2$ & crude size of the statistic
(Lemma~\ref{lem:G-size}(1))\\
$1.6$ & $1 + \Euler < 1.6$ & Euler constant bound\\
$\gamma_{\mathrm{first}}$ & $14.1347$ & rational lower bound for the
first zero height (Input~\ref{lem:zlow}(2))\\
$D_{\mathrm{line}}$ & $\sqrt{251.5^2 + \gamma_{\mathrm{first}}^2}
> 251.89$ & distance from $\Sig$ to any zero with $\beta \le 1/2$\\
$30$ & --- & height below which all zeros are on the line
(Input~\ref{lem:zlow}(3))\\
$17$ & (V4): $h \ge \Gamma/2 + 17$ & slab half-height; covers the
$z$-circle's strip reach $\le 15.9$, plus margin\\
$10^{-6}$ & (V5) & clearance floor for both Cauchy circles\\
$N_b(t)$ & $0.6\log(|t|+2)+2$ & local zero count bound
(Input~\ref{lem:backlund}); $N_{\mathrm{loc}} \le 3$ for
$|t| \le 21.5$\\
$12\log(|t|+4)+12$ & Input~\ref{lem:zeta-prime-prime} & the
non-local part of the $\zeta'/\zeta$ bound\\
$3/2$ & window exponent & architecture constant;
see Remark~\ref{rem:exponent} for the symbolic variant\\
$251.5$ & $\Sig - \tfrac12$ & distance from $z$ to any on-or-left
of-line zero\\
$\dbar, M_1,$ $M_2, D_R$ & \S\ref{sec:derived} & certificate geometry\\
$r_z, \varepsilon_z,$ $W_0, W_1, W_1'$ & \S\ref{sec:derived} & the two
Cauchy circles and their clearances\\
$B_z, B', M_z, M'$ & Definition~\ref{def:Mz} & circle sup-bounds\\
$j_1, j_2, J_0,$ $J_*, U_0, Y_0$ & Definition~\ref{def:thresholds} &
thresholds; $Q(j) \ge C+1$ for $j \ge J_*$\\
$Q(j)$ & $3\lambda^{j+1}/(16M_1M_2) - 52.5$ & the margin function;
$m = Q(j) - C \ge 1$\\
$52.5$ & $C_1/2$ & additive constant in $Q$\\
$L_G(x)$ & $\tfrac12(\log x)^2 + 2\log x + 1.6$ & Lipschitz constant
of $G$ (Lemma~\ref{lem:G-size}(2))\\
$L_h, \delta_{\mathrm{grid}},$ $N_{\mathrm{grid}}$ &
Prop.~\ref{prop:search}; also the precision
$w$ and the error constant $E_0$ & the grid search and its precision\\
\hline
\end{longtable}
}

\appendix

\section{The displayed arithmetic}\label{sec:app:elementary}

All inequalities below are routine verifications from the exponential
series $e^x = \sum_{k \ge 0} x^k/k!$ with its alternating remainder
bound, and from integer arithmetic. They are listed so that a referee
(or the effective shadow) can check each constant of the paper in one
place. (A scripted verification is planned for
\texttt{classical/numerics/}.)

\begin{enumerate}
\item $e^{0.175} < 1.19125$, so $b(0.825) = 0.825\,e^{0.175} <
0.98279 < 0.983 = \eta_*$. \quad
$e^{-0.2} < 0.8188$, so $b(1.2) < 1.2 \cdot 0.8188 = 0.98256 <
\eta_*$.
\item $e^{0.12} > 1.1274$ and $e^{-0.12} > 0.8869$, so
$b(0.88) > 0.88 \cdot 1.1274 = 0.99211 \ge 0.9921$ and
$b(1.12) > 1.12 \cdot 0.8869 = 0.99333 \ge 0.9921$ on
$[v_-,v_+]$.
\item $0.983 \cdot e^{0.008} < 0.9921$; equivalently
$\eta_*/0.9921 < e^{-c}$. The outside thus loses to the inside at net
exponential rate $> c$ (Lemma~\ref{lem:tail}(3)).
\item $e^{0.008} > 1.008$, so $250\,e^{0.008} > 252.00 > 252 \ge
\Sig - \theta$ for all $\theta \in (0,1)$: \eqref{eq:sigma-c}.
\item $251.5^2 = 63252.25$; $\gamma_{\mathrm{first}}^2 < 200$; so
$D_{\mathrm{line}}^2 < 63452.25 < 251.90^2$, and
$D_{\mathrm{line}} > 251.89$.
\item $251^2 + 17^2 = 63290 > 251.5^2$: off-slab off-line zeros are
farther from $z$ than the line is
(Lemma~\ref{lem:circles}(1)).
\item $251^2 + 30^2 = 63901 > 251.90^2 > D_{\mathrm{line}}^2$:
off-line zeros above height $30$ are farther from $\Sig$ than
$D_{\mathrm{line}}$ (Lemma~\ref{lem:circles}(2)).
\item $\sqrt{251.5^2 - 251^2} = \sqrt{251.25} < 15.9$: the
$z$-circle's reach into the strip, whence (V4)'s constant $17$.
\item $\sqrt{251.9^2 - 251^2} < 21.5$: the $\Sig$-circle's reach,
whence $N_{\mathrm{loc}} \le 3$ there.
\item $1.23 \cdot 250 = 307.5$; the rounding margins
$(1.25 - 1.23) \cdot 250\,U \ge 5U - \tfrac12 > 0$ and
$(1.23 - 1.2122)\cdot 250\,U \ge 4.4U - \tfrac12 > 0$ hold for
$U \ge 3$ (Lemma~\ref{lem:j-range}).
\item $52 \cdot e^{0.008} < 52.5$; the bookkeeping of
Lemma~\ref{lem:onesided} then gives \eqref{eq:onesided} with
$C_1 = 105$.
\item $1 + \Euler < 1.6$; $|G(x)| \le 4x^2$ for $x \ge 1$
(Lemma~\ref{lem:G-size}(1), three elementary cases).
\item The function $4x - 3.1(\log x)^2$ is increasing and positive
for $x \ge e$ (used in Lemma~\ref{lem:G-size}(1)).
\end{enumerate}

\begin{thebibliography}{99}

\bibitem{Alp26} L.~Alp\"oge, \emph{Announcement of a Lean-formalized
proof of the Riemann hypothesis}, November 20, 2026; remark at
\texttt{https://alpo.ge/remark.pdf}. The archive was independently
machine-verified on November 21, 2026 \cite{Rub26v}; the present paper
is independent of that outcome.

\bibitem{Rub26v} M.~O.~Rubinstein, \emph{GRH Lean verification},
November 21, 2026,
\texttt{https://mrubinst.github.io/grh-lean-verification}: a public
end-to-end compile of the Alp\"oge archive (all $175{,}187$ modules),
an axiom audit of the five headline theorems, a full kernel replay,
and a check of the $1{,}484{,}729$ declarations they depend on by the
independently written kernel \texttt{nanoda}; scripts and a dated log
included. Further independent verifications by A.~Sutherland (MIT),
E.~Arledge, and S.~DeDeo; tracked on B.~Naskr\k{e}cki's
\texttt{riemann-rev} dashboard.

\bibitem{Bre79} R.~P.~Brent, \emph{On the zeros of the Riemann zeta
function in the critical strip}, Math.\ Comp.\ 33 (1979), 1361--1372.

\bibitem{Edw74} H.~M.~Edwards, \emph{Riemann's Zeta Function},
Academic Press, 1974 (Backlund's explicit zero-counting estimates,
\S1.18).

\bibitem{Ing32} A.~E.~Ingham, \emph{The Distribution of Prime
Numbers}, Cambridge Tracts in Math.\ 30, 1932 (the
P\'olya--Ingham high-derivative method).

\bibitem{Lan05} E.~Landau, \emph{\"Uber einen Satz von Tschebyschef},
Math.\ Ann.\ 61 (1905), 527--550 (Mellin transforms of one-signed
functions; the ``ceiling'' theorem).

\bibitem{Leh66} R.~S.~Lehman, \emph{On the difference $\pi(x) -
\operatorname{li}(x)$}, Acta Arith.\ 11 (1966), 397--410 (the first
zeros: $\gamma_1 = 14.134725\ldots$, on the line).

\bibitem{MV07} H.~L.~Montgomery and R.~C.~Vaughan, \emph{
Multiplicative Number Theory I: Classical Theory}, Cambridge Univ.\
Press, 2007 (\S15.1: Landau's theorem).

\bibitem{PT21} D.~J.~Platt and T.~S.~Trudgian, \emph{The Riemann
hypothesis is true up to $3\cdot 10^{12}$}, Bull.\ London Math.\ Soc.\ 53
(2021), 792--797.

\bibitem{RieVideo} harttraveller, \emph{The Riemann hypothesis,
explained for laypeople} --- video and guide,
\texttt{https://github.com/harttraveller/riemann}; the present paper
formalizes the narration segment ``Step two: the ceiling''.

\bibitem{Rub26} M.~O.~Rubinstein (with AI co-authors), \emph{The
Alp\"oge--Claude proof of the Riemann hypothesis explained for the
Riemann zeta function}, working draft, November 21, 2026,
\texttt{https://www.math.uwaterloo.ca/\textasciitilde
mrubinst/alpoge\_claude\_grh/zeta-paper.pdf}. The windowed Landau
lemma (Part~I \S1.1, Part~II \S1) and its one-page proof are due to
this draft; the present paper effectivizes it.

\bibitem{Tit51} E.~C.~Titchmarsh, \emph{The Theory of the Riemann
Zeta-Function}, 2nd ed.\ (revised by D.~R.~Heath-Brown), Oxford, 1986
(\S\S2.4, 2.12: the Hadamard product and partial fractions for
$\zeta'/\zeta$; \S3.5: the local $\zeta'/\zeta$ representation;
\S9: zero-counting).

\bibitem{Tru14} T.~S.~Trudgian, \emph{An explicit upper bound for the
argument of the zeta function on the critical line}, Math.\ Comp.\ 81
(2012), 1053--1061; and \emph{Updating the zero-counting estimates for
the Riemann zeta function}, J.\ Number Theory.\ 134 (2014), 363--366
(explicit $N(T)$).

\bibitem{Tur53} A.~M.~Turing, \emph{Some calculations of the Riemann
zeta-function}, Proc.\ London Math.\ Soc.\ (3) 3 (1953), 99--117.

\bibitem{Bac14} R.~J.~Backlund, \"Uber die Nullstellen der
Riemannschen Zetafunktion, Acta Math.\ 41 (1918), 345--375 (explicit
zero counting).

\end{thebibliography}

\end{document}
